\documentclass[a4paper,fleqn]{cas-sc}

\usepackage[authoryear]{natbib}

\usepackage{amsthm}

\usepackage{booktabs}
\usepackage{array}
\usepackage{tabularx}
\usepackage{multirow}
\usepackage{threeparttable}
\usepackage{longtable}
\usepackage{xcolor}
\usepackage[normalem]{ulem}
\usepackage{subcaption}
\usepackage{xcolor}
\usepackage[utf8]{inputenc} \usepackage[T1]{fontenc}
\input{lp/lp_numbers}  % LP backtest numbers (generated by paper_tables.py)

\newcolumntype{Y}{>{\raggedright\arraybackslash}X}
\newcommand{\code}[1]{\texttt{#1}}


\newcommand{\best}[1]{\textcolor{red}{\textbf{#1}}}
\newcommand{\secondbest}[1]{\textcolor{blue}{\underline{#1}}}

% More flexible float placement
\setcounter{topnumber}{3}
\setcounter{bottomnumber}{2}
\setcounter{totalnumber}{5}

\renewcommand{\topfraction}{0.90}
\renewcommand{\bottomfraction}{0.80}
\renewcommand{\textfraction}{0.10}
\renewcommand{\floatpagefraction}{0.75}

\begin{document}
\let\WriteBookmarks\relax

\shorttitle{Onchain Monitoring of Stablecoin Depeg Risk}
\shortauthors{O. Chaffard et~al.}
%\begin{frontmatter}

\title[mode=title]{Onchain Monitoring of Stablecoin Depeg Risk}    

% ---------------------------------------------------------------
% Authors and affiliations
% ---------------------------------------------------------------

\author[1,2,7]{Owen Chaffard}[style=chinese,
                    %role=Researcher,
                        orcid=0009-0005-5509-8652]
\cormark[1]
\ead{owen.chaffard@cardoai.com}

\credit{Conceptualization of this study, Data curation, Methodology, Software, etc.}


\author[4,5,7]{Ștefan G\u{a}man}[style=chinese]
\ead{gamanstefan18@stud.ase.ro}
\credit{Conceptualization of this study, Data curation, Background, Methodology, Software}

\author[6,7]{Megang Nkamga Junile Staures}[style=chinese]
\ead{junmeg@ktu.lt}
\credit{Conceptualization of this study, Methodology}
  
\author[1,4,7]{Andrei-Theodor Ginavar}[style=chinese]
\ead{ginavarandrei19@stud.ase.ro}
\credit{Conceptualization of this study, Methodology}

\author[4,7]{Siang-Li Jheng}[style=chinese, orcid=0009-0006-2379-4212]
\ead{jhengsiang24@stud.ase.ro}
\credit{Conceptualization of this study, Data curation, Methodology}

\author[4,7,9]{Daniel Traian Pele}[style=chinese]
\ead{danpele@ase.ro}
\credit{Conceptualization of this study, Supervision, Review \& editing}

\author[3,7,8]{Wolfgang Karl H\"ardle}[style=chinese]
\ead{haerdle@hu-berlin.de}
\credit{Conceptualization of this study, Supervision, Review \& editing}



\affiliation[1]{organization={CardoAI},
            city={Milan},
            country={Italy}}
 
\affiliation[2]{organization={RPTU},
            city={Kaiserslautern},
            country={Germany}}
 \affiliation[3]{organization={School of Data Science, CUEB}, 
        city={Beijing},
        country={China}}

\affiliation[4]{organization={Bucharest University of Economic Studies},
            city={Bucharest},
            country={Romania}}
 
\affiliation[5]{organization={Pozna\'n University of Economics and Business},
            city={Pozna\'n},
            country={Poland}}
\affiliation[6]{organization={Kaunas University of Technology},
            city={Kaunas},
            country={Lithuania}}
\affiliation[7]{organization={IDA Institute Digital Assets},
                city={Bucharest},
                country={Romania}}

\affiliation[8]{organization={School of Business and Economics, Humboldt University of Berlin}, 
        city={Berlin},
        country={Germany}}
\affiliation[9]{Institute for Economic Forecasting, Romanian Academy, Bucharest, Romania}

\cortext[cor1]{Corresponding author}


\begin{abstract}
Stablecoin depegs are often monitored through prices, yet the same price can coincide with materially different executable liquidity. This paper develops a microstructure-aware framework for monitoring market-defined depeg risk in the USDC-USDT Uniswap V3 pool. We construct an hourly dataset of 38,045 observations from January 2022 to May 2026, combining pool states and flows, lending and liquidation activity, oracle prices, liquidity ownership, and market variables. We characterize the liquidity curve through depth, asymmetry, swap-size impact, and an interpretable Gegenbauer decomposition. These inputs support two complementary approaches: an early-warning policy for 15-basis-point depeg events and a probabilistic model forecasting peg-deviation distributions over a 24-hour horizon. The selected early-warning policy detects all five eligible events in the holdout period beginning January 2025, with a median lead time of nine hours and approximately 2.5 false-alert episodes per month. The probabilistic model improves CRPS by 63\% and threshold-weighted CRPS by 37\% relative to a GARCH baseline, with statistically significant gains. In a liquidity-provision backtest on the same pool, setting a vault's range from the forecast intervals raises its risk-adjusted return relative to a Gamma-style vault, and pausing repositioning during early-warning alerts more than halves its losses around depeg events; together, the two models reduce the vault's maximum drawdown by \lpFcKvsGammaMaxDDCut\%. Interpretability and ablation analyses support the predictive relevance of liquidity conditions and broader market variables. Results support liquidity-curve information for venue-level risk monitoring in decentralized markets.
\end{abstract}

\begin{graphicalabstract}
\end{graphicalabstract}

\begin{highlights}
\item Our two-model framework monitors local depeg risk in the USDC-USDT Uniswap V3 pool.
\item Gegenbauer decomposition captures liquidity-curve shape around parity.
\item Tail-weighted CRPS is 14.2\% lower than the strongest benchmark, TiDE.
\item Forecast ranges and early-warning pauses cut LP vault drawdowns in a backtest.
\end{highlights}

\begin{keywords}
Stablecoins \sep Decentralized Finance \sep Early warning
\end{keywords}


\maketitle


\section{Introduction} \label{sec:intro}

Stablecoins have become core infrastructure in crypto markets, where they serve as trading pairs, settlement assets, and collateral. Their importance is reflected in both their use on trading platforms and their market capitalization: the European Central Bank reports that stablecoins were involved in around 80\% of trades on crypto-asset trading platforms in 2025. Stablecoins use different mechanisms to maintain their pegs, and comparative studies identify reserve-backed designs as the most stable and trusted \citep{Yujin2024}. In this context, USDT and USDC stand out as large-cap, long-standing, USD-pegged, reserve-backed stablecoins. Together, they represented 92\% of the market capitalization of major dollar-denominated stablecoins in May 2025 \citep{ECB2025}, and their characteristics make them relatively low-risk assets. Lower risk, however, does not mean no risk: peg stability also depends on redemption access, arbitrage, market liquidity, and confidence in the issuer and its reserves. Reserve backing alone therefore cannot rule out secondary-market price deviations \citep{Lyons2023}. Depeg events, even transient, transmit losses through DeFi markets. Where lending protocols or margin venues mark positions to market, a depeg can reduce collateral values or increase effective leverage, pushing accounts below liquidation thresholds \citep{Perez2021}. Forced sales, liquidation penalties, and slippage can crystallize losses even if the peg later recovers. In automated market maker (AMM) pools, arbitrage against the pool’s stale price can leave liquidity providers holding more of the depegged token and expose them to losses relative to rebalancing \citep{Milionis2022}.

Stablecoin research has examined several related dimensions of this risk \citep{mahrous_2025_stablecoins}. Studies analyze how arbitrage between primary issuance and secondary trading helps stabilize prices \citep{Lyons2023} and how reserve quality and disclosure affect run incentives. Predictive studies test whether market and sentiment variables can predict depegging \citep{stablecoin_depeg_prediction}. Work focused directly on decentralized exchanges has developed pool metrics and changepoint alerts for detecting depegs on Curve StableSwap pools \citep{cintra2023}. Together, these studies show that peg deviations reflect both issuer-level fundamentals and secondary-market conditions. They also motivate a distinction central to this paper: a stablecoin does not have one venue-independent transaction price. The quote observed in an individual pool is formed in that pool's secondary market and need not equal either the issuer's redemption price or a price aggregated across exchanges. Primary- and secondary-market access, as well as the unique trading mechanisms used by centralized exchanges and Automated Market Makers (AMMs), shape how these prices respond to stress \citep{Watsky2024}.

This venue-level perspective matters because traditional AMMs translate pool state into trade prices through an invariant, while Uniswap V3 allows liquidity providers to concentrate capital within chosen price ranges \citep{uniswap_whitepaper}. As a result, the same spot price can coexist with substantially different amounts and distributions of executable liquidity. Liquidity may be asymmetric around the prevailing price, concentrated near parity, or sparse just beyond it. Such differences affect the price impact of order flow and the pool's capacity to absorb trades. Evidence from Uniswap V3 also shows that pool prices can become inaccurate during abrupt stablecoin price movements, underscoring the value of monitoring the local trading mechanism alongside the observed price \citep{Heimbach2022}. On-chain records expose pool states, swaps, and liquidity positions at high frequency, making it possible to study these conditions directly. The open question addressed here is whether detailed, time-varying liquidity conditions can improve venue-level monitoring beyond price and conventional market features.

This paper develops a monitoring framework for a deliberately bounded problem: local deviations from parity in the main USDC--USDT Uniswap V3 pool. The target is the pool's market-defined price, not issuer solvency or the fundamental ability to redeem either stablecoin.

The paper makes two contributions. First, we assemble a detailed empirical dataset from public blockchain records, with 38,045 hourly observations from January 2022 to May 2026. It combines Uniswap V3 and Curve pool states and flows with Aave lending and liquidation indicators, Chainlink prices, liquidity ownership data, and market sentiment. We use these data to construct microstructure features that describe liquidity conditions in greater detail than price and aggregate pool depth alone. In particular, we reconstruct the time-varying Uniswap V3 liquidity curve around both parity and the prevailing price, measure directional swap-size impact, and summarize curve shape with compact Gegenbauer polynomial scores. Together with pool imbalance and order flow, these features represent the level, asymmetry, and local structure of executable
liquidity.

Second, we use this dataset to develop a venue-level risk-monitoring framework with two complementary components: an event-level early-warning policy and a probabilistic model forecasting peg-deviation distributions over a 24-hour horizon. The early-warning component combines classifier scores with Neyman--Pearson threshold calibration and a 12-hour quiet-period alert reset. Model selection prioritizes timely event recall subject to an operational constraint on false-alert burden. Ablations of economically relevant feature blocks assess the contribution of the proposed depeg channels, highlighting in particular the predictive value of Uniswap V3 liquidity-curve shape and Curve 3pool liquidity conditions. In the holdout period beginning January 2025, the selected specification detects all five eligible depeg events, with a median lead time of nine hours and approximately 2.5 false-alert episodes per month.

The second model forecasts the full conditional distribution of peg deviations over the same horizon. Its threshold-weighted scoring objective and flexible distribution, with a spline body and generalized Pareto tails, emphasize large deviations. Exceedance probabilities, quantiles, and expected shortfall provide direct summaries of the likelihood and potential severity of tail outcomes. In the held-out period, this model significantly improves on benchmark forecasters in both overall and tail-focused distributional scores. Together, the models are designed to help investors anticipate periods of heightened pool volatility and assess the observed liquidity and market conditions associated with elevated risk: the classifier signals whether a threshold exceedance event is approaching, while the forecaster characterizes its direction and severity.

Detection rates and scoring rules do not by themselves show that the models are economically useful,
particularly with few held-out events. We therefore evaluate the framework in the decision problem most
exposed to the pool's depeg risk: providing liquidity to it. A liquidity provider earns swap fees but loses
to arbitrageurs whenever the price moves, a cost measured by loss-versus-rebalancing (LVR)
\citep{Milionis2022}, and automated vaults that recentre a narrow range on the price amplify that cost
during depegs. In a backtest against a Gamma-style vault, each model has one role. The forecaster sizes the
range: using its predictive interval raises the vault's net profit and Sharpe ratio and reduces its
repositioning costs and daily tail losses. The classifier pauses repositioning during its alerts: this
cuts the LVR incurred around depeg events by \lpFcKEvLVRCut\% and the drawdowns within them by
\lpFcKEventDDCut\%, more than a price trigger at the same threshold or randomly timed pauses achieve. The
pause gives up part of the high fee income of depeg periods, so it acts as insurance rather than as a source
of return. A stress scenario based on the March 2023 USDC depeg, in which vaults that keep repositioning lose
20--30\% of their capital, illustrates the losses it protects against.

The remainder of the paper is organized as follows. Section \ref{sec:background} describes the institutional setting and the on-chain channels relevant to the stated scope. Section \ref{sec:methodology} presents the data, liquidity-curve representation, and modeling approach. Section~\ref{sec:results} reports the empirical results, including the economic evaluation in Section~\ref{sec:lp_results}, and Section~\ref{sec:conclusion} concludes.


\section{Background and institutional setting}
\label{sec:background}
\subsection{Scope and Definitions}

The scope of this paper is not general prediction of stablecoin depeg risk, but the development of an on-chain framework for continuous risk monitoring. We apply this framework empirically to provide early warning indicators of potential market dislocations in the USDC-USDT Uniswap V3 pool, a major stablecoin trading venue. The design of the framework and the choice of empirical setting are guided by three central distinctions in classifying depeg risk.

First, it is essential to distinguish the \emph{fundamental} peg from the \emph{market-defined} peg. The fundamental peg is anchored by the ability to redeem, mint, or otherwise convert the stablecoin at or near its reference price \citep{Lyons2023}. The market-defined peg is the price at which the stablecoin actually trades on a given venue and is therefore local and highly dependent on market microstructure \citep{MAKAROV2020} . As a result, a stablecoin can experience a transient local market depeg even when its fundamental peg remains intact, as illustrated by the USDe dislocation on Binance on October 10, 2025, when the stablecoin briefly traded as low as \$0.65 on that platform due to an internal oracle design issue, while prices remained stable on major onchain venues.

Second, depeg risks can be \emph{exogenous} or \emph{endogenous}. Exogenous depegs originate outside the trading venue and include regulatory and backing-related risk channels. Such events typically arise as discrete shocks and are difficult to anticipate before the relevant information becomes public. Endogenous depegs, such as those driven by sentiment or liquidity risk, arise within the local market structure and are directly reflected in on-chain pool states and transaction flows.

Third, depegs may be transient or persistent \citep{DUAN2023}. A transient depeg is a short-lived deviation that resolves as arbitrageurs rebalance across venues or liquidity returns, whereas a persistent depeg reflects a more durable impairment of confidence, convertibility, or market depth. The collapse of the Terra–Luna ecosystem in May 2022 exemplifies a persistent loss of peg and has been extensively studied \citep{BRIOLA2023}. Transient events are not necessarily benign from a risk-management perspective: even brief price deviations in DeFi can cause irreversible losses through liquidations, adverse execution, or loss versus rebalancing in narrow liquidity provision \citep{QIN2021}.

These distinctions define the object of this study. We do not attempt to forecast rare issuer-side failures or regulatory shocks. Instead, we focus on local, endogenous, and potentially transient depeg risk in the USDC--USDT market. This is the class of risk for which active monitoring and hedging are most feasible and for which an on-chain approach is most informative. Public blockchains provide real-time granular information on market microstructure, liquidity conditions, order flow, and interactions between venues, allowing one to observe how stress is absorbed, amplified, or transmitted between linked DeFi venues. For stablecoin depeg risk, this transparency is particularly valuable: relevant warning signals often manifest not in the price level itself, but in the evolving distribution of liquidity around the peg, persistent inventory imbalances, and the breakdown of arbitrage across venues.

\subsection{Depeg Channels and On-Chain Observability}
\label{subsec:depeg_channels_observability}

Beyond the distinctions introduced above, depeg risks can be organized into several economically meaningful channels. However, for an early-warning framework, the key dimension is not only \emph{why} a peg deviation occurs, but also whether the underlying stress generates observable signals before a price dislocation is realized. Two primary risk channels concern backing and convertibility: (i) issuer reserve concerns or redemption impairments, and (ii) regulatory restrictions and access frictions \citep{eichengreen2025}. Both channels are exogenous to the trading venue, directly affect the fundamental peg, and typically arrive as discrete information events with limited scope for advance detection from onchain observables. Given the focus of this study, these channels are largely outside our empirical remit.

Our framework instead concentrates on two predominantly endogenous and on-chain observable channels. The first is liquidity and microstructure risk, which constitutes the main object of analysis in this paper. On-chain markets feature idiosyncratic liquidity conditions and trading microstructure, while public blockchain indexes provide detailed high-frequency data on pool states, liquidity depth and ownership, as well as order flow composition; these inputs form the core feature block of our dataset. 

The second channel comprises sentiment and run dynamics, which are also substantially observable on-chain. This channel motivates our focus beyond the primary venues of stablecoin trading toward cross-protocol amplification mechanisms, including lending pressures, liquidation cascades, and global market sentiment proxies. By incorporating these additional dimensions, the framework captures how localized stress in a stablecoin pool can propagate through the broader DeFi ecosystem.

The remainder of this background section therefore focuses on the DeFi microstructure through which these observable channels operate. In particular, the next subsection shows how Curve and Uniswap V3 provide complementary information about stablecoin stress: Curve captures broad stablecoin-pool imbalances, whereas Uniswap V3 reveals the endogenous shape of executable liquidity around the USDC--USDT peg.

\subsection{AMM Microstructure and Liquidity-Driven Depeg Severity}

\subsubsection{AMM liquidity and local price formation}

\begin{figure}[pos=!htbp,align=\centering]
    \centering
    \includegraphics[width=0.5\linewidth]{figures/invariant_with_slippage.png} 
    \caption{Univariate curve of a constant product AMM. An example of a swap is provided which highlights the large price impact (slippage) of transactions \href{https://github.com/QuantLet/onchain-insights/tree/main/2.\%20Curve\%20liquidity\%20pools}{\includegraphics[height=0.6\baselineskip]{figures/Quantlet LOGO.png}}}
    \label{fig:AMM}
\end{figure}

DeFi trading microstructure is organized around automated market makers (AMMs), which implement on-chain trading mechanisms that replace traditional limit-order books with deterministic pricing rules encoded in smart contracts \citep{hypeorhope}. Liquidity is supplied by users, or liquidity providers (LP), who deposit token pairs or baskets into pool contracts, against which trades are executed as quoted prices adjust mechanically with changes in pool reserves. The resulting liquidity depth determines slippage, that is, the mechanical price adjustment induced by a given trade. Traditional AMMs employ a constant-product invariant that guarantees non-exhaustible liquidity across all possible prices, but at the cost of substantial price impact even for relatively small orders (Figure~\ref{fig:AMM}). Stablecoin-focused AMMs instead concentrate liquidity in a narrow band around the expected peg, thereby increasing effective depth near parity. This design choice suppresses volatility in normal conditions by making prices locally inelastic, but may amplify large price dislocations due to liquidity thinning away from parity. The two main protocols in our study implement different mechanisms for concentrating liquidity, which in turn generate distinct but complementary predictive signals for stablecoin price formation.

\subsubsection{Curve Finance: Pool Imbalance as Broad Stablecoin Stress}

Curve Finance's StableSwap invariant extends the standard AMM design to pools containing multiple closely priced assets, with an emphasis on stablecoins. Its principal example is the 3pool, which comprises USDC, USDT, and DAI. The invariant creates a parameterized flat pricing region around parity, governed by an amplification parameter $A$ set through DAO governance, thereby concentrating liquidity symmetrically near the target exchange rate (Figure~\ref{fig:liqconcentration1}). 

This design enables large stablecoin trades at low slippage when the pool is balanced. During periods of stress, however, the pool may accommodate persistent one-sided flows through substantial inventory imbalances before pronounced price deviations emerge. Consequently, early-warning indicators should extend beyond prices and trading volumes to include pool composition, concentration or entropy measures, and marginal slippage.

\subsubsection{Uniswap V3: Endogenous Liquidity Shape}

Uniswap V3 reduces slippage through \emph{concentrated liquidity}, a mechanism that differs from Curve's StableSwap design. Rather than modifying the AMM invariant, Uniswap V3 retains the constant-product structure while allowing liquidity providers (LP) to allocate capital within self-selected price intervals. This feature improves capital efficiency relative to conventional AMMs, in which liquidity is distributed across the entire price range. 

Price intervals are discretised into ticks, with one tick corresponding to one basis point; at each tick, active liquidity determines the order flow that can be absorbed before the price moves to the next interval. Active liquidity equals the aggregate liquidity of all LP positions whose selected ranges include the underlying tick.

The resulting active-liquidity curve, defined over the tick grid, aggregates heterogeneous LP allocation strategies into functional, endogenous and time varying liquidity concentration (figure \ref{fig:liqconcentration2}). Since fee accrual requires LP positions to remain within the current price range the curve is directly linked to LP strategies. Its functional shape may exhibit asymmetries, heavy tails, fragmentation, or liquidity cliffs that are not captured by scalar measures of pool imbalance or aggregate depth. Consequently, pools with similar near-peg liquidity can differ substantially in resilience depending on whether available liquidity is balanced, skewed, dispersed, or concentrated within narrow price ranges. These considerations motivate the functional analysis of the USDC--USDT liquidity curve in our approach (Section \ref{sec:liq_curve_decomposition}).

\begin{figure}[pos=!htbp,align=\centering]
    \centering
    \begin{subfigure}[b]{0.475 \textwidth}
        \centering
        \includegraphics[width=\linewidth]{figures/curve_liquidity_paper.png}
        \caption{Stableswap liquidity distribution \label{fig:liqconcentration1}}  
    \end{subfigure}
    \hfill
    \begin{subfigure}[b]{0.49\textwidth}
        \centering
        \includegraphics[width=\linewidth]{figures/owner_zones.png}
        \caption{Uniswap v3 liquidity distribution \label{fig:liqconcentration2}}    
    \end{subfigure}
\caption{Liquidity concentration in Stableswap and Uniswap v3 protocols. The Stableswap liquidity distribution is fixed, with its shape controlled by a constant parameter (which can only be changed by a DAO decision). On the other hand, the Uniswap v3 liquidity curve is dynamic, and the result of overlapping LP positions on fixed price segments. The 10 largest positions and their owners (truncated wallet addresses) are highlighted.\href{https://github.com/QuantLet/onchain-insights/tree/main/3.\%20Uniswap\%20liquidity\%20curve}{\includegraphics[height=0.6\baselineskip]{figures/Quantlet LOGO.png}}}
\label{fig:liqconcentration}
\end{figure}

\subsection{Cross-protocol amplification channels}
\label{sec:ampl_channels}

The AMM microstructure described above captures the most direct on-chain manifestation of stablecoin stress: price deviations, pool imbalances, slippage, and the distribution of liquidity around the peg. However, depeg risk is not only a spot-market phenomenon. Stress can also emerge or be amplified through the ownership structure of liquidity, lending markets, oracle-based liquidation mechanisms, and broader changes in market sentiment. These additional channels motivate the remaining components of our empirical dataset.

\paragraph{Liquidity ownership and effective depth}

A central challenge in assessing on-chain liquidity is that displayed liquidity may overstate the capital effectively available to absorb market stress. This distinction is particularly relevant in Uniswap V3, where liquidity providers can dynamically adjust or withdraw positions across selected price ranges. Maximum extractable value (MEV) activity, particularly just-in-time (JIT) liquidity provision, further complicates this assessment \citep{xiong2023}. JIT providers may condition liquidity provision on anticipated uninformed order flow \citep{capponi2024}, causing observed liquidity to overstate the depth available during adverse market conditions. For early-warning purposes, aggregate liquidity measures should therefore be complemented by indicators of liquidity ownership, concentration, and persistence.

\paragraph{Aave lending markets}

Stablecoin stress may also propagate through decentralised lending markets. USDC and USDT are widely used as collateral on Aave, where funding conditions and liquidation mechanisms can amplify on-chain peg instability. Rising utilisation and borrowing rates may reflect funding pressure, while collateral liquidations can generate feedback effects that reinforce price deviations, run dynamics, and market imbalances \citep{P2Plending}. Accordingly, utilisation, borrowing rates, and aggregate liquidation activity constitute relevant early-warning indicators.

\paragraph{Oracle feeds and broader market state}

Oracle prices constitute a further transmission channel for stablecoin stress. DeFi protocols do not use decentralised-exchange prices directly for collateral valuation and liquidation decisions, but instead rely on external reference feeds. These feeds incorporate broader market conditions, including major cryptocurrency prices and the stability of the underlying fiat peg, and may therefore affect whether observed liquidity stress develops into a realised price dislocation. We consequently include the volatility of major cryptoassets, notably ETH and BTC, as well as USD volatility and a crypto-sentiment index as proxies for prevailing market conditions.
\paragraph{}
This institutional setting motivates the empirical design of the paper. We construct a dataset that combines AMM microstructure variables, lending-market indicators, oracle reference prices, and market-wide sentiment measures (full feature set available in Appendix~\ref{sec:data-appendix}). The following section describes the collection, synchronization, and feature construction procedures used to translate these channels into inputs for the early-warning model.

\section{Methodology}\label{sec:methodology}

 \subsection{Data collection} \label{sec:data_collection}

Our early warning models are constructed from a comprehensive dataset of on-chain conditions for the USDC--USDT pair and surrounding markets. We collect hourly data from each of the previously discussed protocols through their publicly indexed subgraphs on The Graph \citep{grtwhitepaper}. Additional information on liquidity ownership is obtained through the public Etherscan API, while observations on market sentiment are collected from the publicly available crypto fear and greed index API (details on data sources are available in appendix~\ref{sec:data-appendix}).

The sample used in this study spans January 2022 to May 2026 and contains 38{,}045 hourly observations; for reproducibility, the dataset is updated daily and made publicly available in its latest version\footnote{\url{https://doi.org/10.5281/zenodo.18339067}} \citep{Defidataset}. All observations are indexed hourly and aligned to the nearest finalized Ethereum block at each hour to avoid temporal inconsistencies across venues. Protocol state variables are queried at the corresponding block, while flow-based measures such as volume and order flow are aggregated over the interval between consecutive hourly blocks in the common time index.

We define the endogenous price of the USDC--USDT pair as the price in the Uniswap v3 USDC-USDT pool with a 0.01\% fee tier, measured from the pool state at each hourly observation. By convention, this price is quoted as the price of USDC in units of USDT. This pool is the most actively traded venue for the pair due to its low fee and offers the highest price resolution, with one tick corresponding to one basis point. Accordingly, all subsequent references to the Uniswap liquidity curve pertain to the liquidity distribution in this pool. We collect hourly snapshots of the active liquidity curve within a $\pm 50$-tick window around the peg and, separately, around the current active tick in order to derive swap-impact curves and identify liquidity cliffs (see Section~\ref{sec:data_processing}).

All protocols considered in the analysis were active at the chosen start of the sample period, so the resulting dataset contains no missing values arising from protocol availability. Our data pipeline accounts for the transition of the Aave protocol from v2 to v3 by selecting, at each hourly observation, the market with the larger total value locked. By contrast, we do not incorporate the transition to Uniswap v4, as the v3 USDC-USDT pool remains the dominant venue for this pair and continues to be the preferred venue for liquidity provision and trading. Additional details on the dataset, including the underlying data sources and the full list of features used in the analysis, are provided in Appendix~\ref{sec:data-appendix}.


\subsection{Data preprocessing}
\label{sec:data_processing}

The objective of our data preprocessing is to transform heterogeneous, granular on-chain observations into a compact representation of each pool’s capacity to absorb directional trading pressure around parity, aligned with the risk channels outlined in Section~\ref{sec:background}.

\subsubsection{Feature blocks motivated by on-chain depeg channels}

\paragraph{Realised price and flow pressure}
The first feature block captures the current pool price and contemporaneous order-flow pressure. The Uniswap v3 pool price is recorded as its deviation from parity, expressed in basis points. Order flow is summarized through both transaction counts and traded volume, together with net (signed) flow into USDC (respectively, USDT) which directly reflects directional pressure on each stablecoin (Figure \ref{fig:swaps1}). In addition, the same volume-based metrics are computed for the corresponding higher fee tier pool, since trading activity in a higher fee tier may signal market expectations of elevated short-term volatility.

\paragraph{Aggregate pool imbalance and effective depth} \label{sec:liq_imbalance}

The second feature block measures accumulated directional imbalance beyond the instantaneous price deviation. The Curve 3pool provides an aggregate view of the relative demand for USDC and USDT, and we therefore include the USDC and USDT balance proportions in this pool. A compact scalar measure of the overall 3pool imbalance is given by the mixing Shannon entropy \citep{cintra2023}:
\begin{equation}
S_{\mathrm{3pool}} = -\left(w_{\mathrm{DAI}}\log w_{\mathrm{DAI}} + w_{\mathrm{USDC}}\log w_{\mathrm{USDC}} + w_{\mathrm{USDT}}\log w_{\mathrm{USDT}}\right),
\end{equation}
where $w_{\mathrm{DAI}} + w_{\mathrm{USDC}} + w_{\mathrm{USDT}} = 1$ denote the normalized token balances in the pool.

Complementing this aggregate measure, the price-centered liquidity curve from the main Uniswap v3 pool provides a more detailed characterization of liquidity asymmetries. To characterize the pool’s ability to absorb trades, we construct directional swap-size impact curves around the current active tick (Figure \ref{fig:swaps2}). For each candidate tick \(k\), we calculate the cumulative active liquidity available between the current active tick \(k_t\) and \(k\). For an upward move, this is the liquidity encountered as a trade pushes the price toward higher ticks; for a downward move, it is the corresponding liquidity encountered toward lower ticks. The resulting functions summarize how much liquidity a trade must traverse before moving the pool price to a given tick:
\begin{equation}
\begin{split}
\operatorname{Swap}_t^{+}(k) &= \sum_{j=k_t}^{k} L_{t,j}, \qquad k>k_t,\\
\operatorname{Swap}_t^{-}(k) &= \sum_{j=k}^{k_t} L_{t,j}, \qquad k<k_t.
\end{split}
\end{equation}
where \(L_{t,j}\) denotes active liquidity in tick interval \(j\). In implementation, these cumulative liquidity measures are converted into the corresponding token amount required to move the price across the tick grid. Thus, the impact curve maps a desired price displacement to the size of the swap required to produce it.

The upward and downward curves provide complementary measures of liquidity resilience. A high curve indicates that substantial trading volume is required to move the price through that region, whereas a low curve indicates that the price can move relatively easily. Discontinuities or sharp declines in the curve identify liquidity cliffs: regions in which a small additional trade can generate a comparatively large price movement.

We summarize the curves at two scales. First, the area under each directional impact curve provides a global measure of effective liquidity imbalance. Second, the tangent of each curve at the prevailing price measures local effective imbalance. It captures the immediate marginal depth available to absorb a small trade, before the trade reaches more distant liquidity bands. A large difference between the upward and downward tangents therefore indicates that the pool is locally more fragile on one side of the peg, even if aggregate liquidity over the wider range appears relatively balanced.

These measures therefore allow us to exploit the endogenous liquidity curve of the Uniswap v3 pool to distinguish broad resilience from near-price fragility, by separately quantifying aggregate depth over a wider stress region and the marginal depth immediately around the prevailing price.


\begin{figure}[pos=!htbp,align=\centering]
    \centering
    \begin{subfigure}[b]{0.48\textwidth}
        \centering
        \includegraphics[width=\linewidth]{figures/Net_swap_hourly_SMA.png}
        \caption{Moving average of the hourly net swap flows\label{fig:swaps1}}
    \end{subfigure}
    \begin{subfigure}[b]{0.48\textwidth}
        \centering
        \includegraphics[width=\linewidth]{figures/swap_size_impact.png}
        \caption{Snapshot of the swap size impact curves \label{fig:swaps2}}
    \end{subfigure}
    \caption{Swap size price impact curve and rolling average of net flows in the pool. When compared to the trend of hourly net swaps, the slippage can be a predictive leading indicator of heightened volatility events which may lead to a loss of peg. The highlighted snapshot shows an imbalance in the pool favoring USDT (smaller USDC buys lead to larger slippage) \href{https://github.com/QuantLet/onchain-insights/tree/main/3.\%20Uniswap\%20liquidity\%20curve}{\includegraphics[height=0.6\baselineskip]{figures/Quantlet LOGO.png}}}
    \label{fig:swaps}
\end{figure}

\paragraph{Broader market and lending conditions}
Local liquidity conditions do not determine depeg risk in isolation. The same imbalance or liquidity cliff may remain contained under benign market conditions but become destabilizing when leveraged positions, collateral constraints, or broader risk aversion amplify trading pressure. We therefore augment the AMM-level features with indicators capturing lending-market conditions, cryptoasset market stress, dollar-market conditions, and aggregate sentiment.
For lending markets, we collect Aave utilization, borrowing rates, and liquidation activity. High utilization and elevated borrowing costs may indicate that liquidity is becoming scarce or that traders are relying more heavily on leverage. Liquidation activity provides a direct measure of realized deleveraging pressure: forced sales can increase order-flow imbalance, consume available AMM liquidity, and reinforce the initial price deviation. These variables therefore act as indicators of whether a local liquidity shock is likely to propagate through collateral and liquidation channels.

We also include trend and volatility measures for BTC and ETH to control for the broader crypto-market environment. Rolling returns, moving averages, and volatility statistics capture changes in market direction, uncertainty, and risk appetite that may affect stablecoin demand and the willingness of liquidity providers to maintain positions. 

To capture conditions in dollar markets, we construct a synthetic USD index from the principal USD exchange-rate pairs available through Chainlink oracles:

\begin{equation}
\log(FX^{USD}_t)
=
0.5\log(p^{EUR}_t)
+
0.3\log(p^{JPY}_t)
+
0.2\log(p^{GBP}_t),
\end{equation}
where each \(p_t\) denotes the USD price expressed in the corresponding currency. We derive rolling changes, trends, and volatility measures from this index. These variables provide a broad proxy for dollar-market conditions and may capture shifts in the external demand for dollar exposure that are not visible from the AMM state alone.
Finally, we include a crypto fear-and-greed index as a summary measure of market-wide sentiment. Sentiment is not interpreted as a direct causal driver of depeg events; rather, it provides a conditioning variable for the response of traders and liquidity providers to local stress. In combination, these features describe the external environment in which AMM fragility operates. The resulting framework distinguishes between isolated microstructural weakness, which may be absorbed through arbitrage without a major price dislocation, and microstructural weakness occurring alongside leverage, liquidation, market-volatility, or risk-aversion pressures, where the probability of propagation is substantially greater.

\subsubsection{Functional decomposition of the endogenous liquidity curve}
\paragraph{The peg-centered functional liquidity representation}
Aggregate liquidity measures describe how much liquidity is available in the pool, but not where that liquidity is located relative to parity. This distinction is central to stablecoin depeg risk. Liquidity concentrated close to the peg can support routine trading while providing limited protection against a larger price displacement. Conversely, liquidity positioned farther from parity determines the pool’s ability to absorb trades once the price begins to move away from its intended value.

We therefore represent the Uniswap v3 liquidity distribution in a coordinate system centered on the stablecoin peg rather than on the contemporaneous active tick. Let \(k^{\mathrm{peg}}\) denote the tick corresponding to parity and \(k\) denote an observed tick. We map the peg-centered tick window to a normalized coordinate \(x\in[-1,1]\), with \(x=0\) corresponding to the peg. At each hour \(t\), the resulting liquidity curve is denoted by \(L_t(x)\). We work with log-liquidity,
\begin{equation}
\ell_t(x)=\log\{L_t(x)+1\},
\end{equation}
The log transformation reduces the influence of very large liquidity values and makes proportional changes in liquidity more comparable across the tick range.

This representation differs from the current-price-centered representation used to construct the swap-size impact curves. The impact curves are naturally evaluated relative to the current active tick because they measure the liquidity encountered by a trade starting from the prevailing price. The lower-dimension representation, by contrast, is intended to track the structural shape of liquidity provision relative to the economically meaningful reference point: parity. A price-centered representation would shift the entire curve whenever the active price moved, potentially making a stable liquidity profile appear to change simply because the coordinate system had moved (see Appendix \ref{sec:liqcurve}). Peg-centering keeps the reference frame fixed and makes the resulting coefficients comparable over time.

The peg-centered curve is especially useful because it separates two distinct sources of variation. Changes in the overall level of the curve indicate broad liquidity withdrawal or replenishment, while changes in its shape indicate concentration, asymmetry, depletion near parity, or the emergence of localized gaps. These shape changes are economically relevant even when the observed price remains close to the peg. In particular, a reduction in liquidity on one side of parity may reveal directional fragility before it produces a large realized price deviation.

After constructing the hourly peg-centered curves, we approximate each curve using a low-dimensional fixed basis of Gegenbauer polynomials.

We choose the Gegenbauer family for two reasons. First, its orthogonality weight, $w_{\alpha}(x)=(1-x^2)^{\alpha-\frac{1}{2}},$ allows $\alpha$ to adjust the relative emphasis assigned to the center and endpoints of the normalized peg-centered window. Because the forecasting objective places greater weight on large deviations, but the appropriate emphasis across the liquidity curve is not known a priori, we treat $\alpha$ as a tunable hyperparameter and select it empirically. This lets the representation adapt its sensitivity to liquidity farther from parity without fixing an arbitrary weighting in advance. Second, the baseline representation should treat both sides of parity symmetrically: USDC and USDT share the same target peg, so neither direction should receive greater weight ex ante. The Gegenbauer weight is even in $x$ and provides this symmetric reference measure. This does not impose symmetry on the realized liquidity curve: odd-degree components can capture directional asymmetries, which are part of the depeg-risk signal. The more general Jacobi family permits unequal weights on the two sides, adding an a priori directional preference that is not part of our baseline assumption.

\paragraph{Interpretable Gegenbauer decomposition}
\label{sec:liq_curve_decomposition}

\begin{figure}[pos=!htbp,align=\centering]
    \centering
    \includegraphics[width=\linewidth]{figures/gegenbauer_curve_reconstruction.png}
    \caption{Observed peg-centered log-liquidity curve and its truncated 7 degree Gegenbauer reconstruction for a representative hourly observation. The middle panel reports the residual approximation error across the normalized tick range, while the lower panel shows the corresponding Gegenbauer coefficients by polynomial degree. These coefficients provide a compact representation of the liquidity curve’s level, asymmetry, curvature, and higher-order localized shape variation.\href{https://github.com/QuantLet/onchain-insights/tree/main/7.\%20Gegenbauer\%20Polynomials}{\includegraphics[height=0.6\baselineskip]{figures/Quantlet LOGO.png}}}
    \label{fig:gegen_decomposition}
\end{figure}

The peg-centered liquidity curve is a function of relative tick position
rather than a single scalar. Before decomposition, we subtract the mean
log-liquidity curve estimated exclusively from the training observations.
Let $\mathcal{T}$ denote the training set and define
\begin{equation}
\bar{\ell}_{\mathcal{T}}(x)
=
\frac{1}{|\mathcal{T}|}
\sum_{s\in\mathcal{T}}\ell_s(x),
\qquad
\widetilde{\ell}_t(x)
=
\ell_t(x)-\bar{\ell}_{\mathcal{T}}(x).
\end{equation}
The training mean curve is held fixed when transforming subsequent
validation and test observations, preventing future observations from
entering the centering operation. This centering concerns liquidity
levels at each tick and is distinct from centering the tick coordinate
around parity.

We approximate the centered curves using a fixed Gegenbauer basis:
\begin{equation}
\widetilde{\ell}_t(x)
\approx
\sum_{k=0}^{K}
\beta_{t,k}\phi_k^{(\alpha)}(x).
\end{equation}
The original log-liquidity curve is reconstructed by adding
$\bar{\ell}_{\mathcal{T}}(x)$ to the truncated expansion. The coefficient \(\beta_{t,k}\) measures how strongly the liquidity curve at time \(t\) resembles the \(k\)-th shape pattern in the basis. In the continuous formulation, the coefficient is obtained by projection,
\begin{equation}
\beta_{t,k}
=
\int_{-1}^{1}
\ell_t(x)\phi_k^{(\alpha)}(x)w_\alpha(x)\,dx,
\end{equation} with the empirical implementation using the corresponding weighted projection over the discrete tick grid. Thus, the decomposition replaces a high-dimensional liquidity snapshot with a compact vector of shape scores.

The coefficients describe departures from the training mean
log-liquidity curve. The degree-zero component captures broad changes
in liquidity level relative to this reference. The degree-one component
captures changes in left-right asymmetry around parity, while the
degree-two component describes changes in broad curvature, including
concentration near parity relative to the outer regions. Higher-degree
components capture finer spatial departures from the reference curve.

The polynomial degree should therefore be understood as a scale of spatial shape complexity rather than as a literal frequency.
This representation is particularly useful for depeg-risk analysis because it distinguishes symmetric liquidity changes from directional ones. For example, broad liquidity withdrawal lowers the degree-zero score,
whereas uneven withdrawal across the two sides of parity also changes
the odd-degree scores. A pool can therefore exhibit little change in total liquidity while becoming substantially more fragile in one direction.
We retain the first \(K+1\) components according to the cumulative proportion of variance explained in the training set by the truncated reconstruction. This provides dimensionality reduction while preserving the dominant patterns of variation in the liquidity curves. The use of a fixed basis also ensures that the scores remain comparable over time. This differs from rolling functional principal component analysis, in which the estimated eigenfunctions can change across windows and make coefficient movements more difficult to interpret.
In addition to the individual coefficients, we construct grouped energy measures that summarize variation across economically related shape scales. For example,
\begin{equation}
\begin{split}
E_t^{\mathrm{odd}}&=\sum_{k\in\{1,3,5,\ldots\}}\beta_{t,k}^{2} \\
E_t^{\mathrm{even}}&=\sum_{k\in\{2,4,6,\ldots\}}\beta_{t,k}^{2}.
\end{split}
\end{equation}

Odd-degree energy summarizes asymmetric deformation around the peg, while even-degree energy summarizes more symmetric changes in concentration and curvature. We also group coefficients into low-, middle-, and high-degree bands. Low-degree energy captures broad changes in the liquidity landscape, whereas high-degree energy captures more localized features that may correspond to liquidity notches or cliffs.

Finally, we calculate rolling means and rolling volatilities of both the individual coefficients and the grouped energy measures. These temporal features distinguish persistent changes in liquidity structure from short-lived distortions. The resulting Gegenbauer feature set therefore captures three dimensions of liquidity fragility: (i) the level of available liquidity, (ii) the shape of its distribution around parity, and (iii) the instability of that shape over time. The parameter \(\alpha\), which controls the relative emphasis placed on the centre and tails of the tick interval, is treated as a model hyperparameter and is discussed in Appendix \ref{sec:liqcurve}.


\subsubsection{Liquidity ownership analysis}

The aggregate liquidity curve describes the amount and location of liquidity available in the pool, but it does not indicate who supplies that liquidity or how persistent it is likely to be. This distinction matters because displayed liquidity may overstate the pool’s effective capacity when a large share is supplied by short-lived, transaction-specific positions. We therefore complement the curve-level features with measures of liquidity ownership and position structure. 

We classify liquidity positions according to the mechanism through which they are created. Positions minted through the Uniswap \texttt{NonfungiblePositionManager} (NFPM) are treated as externally managed positions associated with identifiable liquidity providers. Positions created directly by other smart contracts are classified separately, as these positions are often associated with automated strategies, including just-in-time (JIT) liquidity provision and other MEV-related activity.

We classify non-NFPM liquidity according to the heuristic in \cite{xiong2023}. Liquidity minted and burned within the same ethereum block is thus marked as JIT. We then calculate the share of total minted liquidity associated with these short-lived positions (Figure \ref{fig:jit_share}). This measure captures the extent to which the pool’s apparent liquidity depends on provision that is activated only around individual transactions. A high JIT share may support execution under normal conditions, but it may provide less persistent protection against a sequence of trades or a sustained depeg episode.

For NFPM positions, we construct additional features describing the concentration and structure of liquidity that remains in range (Figure \ref{fig:positions_feats}). These include the Herfindahl-Hirschman Index (HHI), top-\(k\) liquidity shares, the liquidity-weighted average width of positions, and the liquidity-weighted average age of positions. 

The concentration measures identify whether effective liquidity is broadly distributed or dominated by a small number of positions. Narrow positions indicate that liquidity is concentrated over a limited tick range and may therefore be highly sensitive to price movements away from the peg. Position age provides a complementary measure of persistence, distinguishing recently created liquidity from positions that have remained in the pool over a longer period.

The ownership features are therefore included alongside, rather than as substitutes for, the aggregate liquidity measures. Together they capture both the shape of available liquidity and the structural conditions underlying that liquidity, allowing the early-warning models to assess whether observed depth is likely to remain available as market stress develops.

\begin{figure}[pos=!htbp,align=\centering]
    \centering
    \begin{subfigure}[b]{0.465\textwidth}
        \centering
        \includegraphics[width=\linewidth]{figures/jit_share_over_time.png}
        \caption{Share of JIT liquidity among non-NFPM positions \label{fig:jit_share}}
    \end{subfigure}
    \begin{subfigure}[b]{0.495\textwidth}
        \centering
        \includegraphics[width=\linewidth]{figures/hhi_tickwidth_over_time.png}
        \caption{NFPM positions HHI and average price width\label{fig:positions_feats}}
    \end{subfigure}
    \caption{Historical plots of liquidity-position characteristics for the Uniswap USDC-USDT v3 pool. The first plot documents the recent increasing share of pool liquidity attributable to JIT provision, which raises the risk of a systematic overestimation of the pool's effective reserves. The second plot presents key features of traditional NFPM positions, indicating broader structural changes among LPs. The HHI shows that liquidity has become less concentrated across providers, implying a more diffuse ownership structure. At the same time, the average width of liquidity positions has increased, suggesting that the pool reserves are now composed of more numerous, lower-value positions with lower individual risk exposure.\href{https://github.com/QuantLet/onchain-insights/tree/main/4.\%20Stablecoin\%20liquidity\%20ownership}{\includegraphics[height=0.6\baselineskip]{figures/Quantlet LOGO.png}}}
    \label{fig:positions}
\end{figure}


\subsection{Modeling Approach}
\label{sec:method-modeling}

Our data-processing pipeline translates heterogeneous on-chain observations into a structured representation of stablecoin market resilience. The resulting dataset combines realized price and order-flow pressure, aggregate pool imbalance, local swap-price impact, liquidity ownership, lending-market conditions, and broader crypto-market state variables. It also includes a low-dimensional representation of the time-varying Uniswap v3 liquidity curve around the peg. These features capture not only whether the pool is currently under pressure, but also how much directional trading pressure it can absorb and whether its available liquidity is concentrated, asymmetric, or potentially transient.

We use this dataset to develop two complementary predictive models, each designed for a distinct risk-monitoring objective. The first is a binary classifier that estimates whether a material depeg will occur within the subsequent 24-hour warning window. This model is intended to provide an operationally interpretable warning signal.

The second model is a probabilistic forecasting neural network that estimates the conditional distribution of future stablecoin prices rather than only a binary event probability. Its hybrid architecture combines temporal representations of the observed market state with sparse covariate selection and a flexible distributional output. This enables the model to characterize the direction, magnitude, and tail risk of future price deviations, including the probability of exceeding thresholds relevant for collateral and liquidation risk.

The two models therefore use the same microstructure-rich dataset to answer complementary questions. The classifier determines whether intervention may be required within a specified warning horizon, while the probabilistic model quantifies the potential severity and distribution of the ensuing price movement. Together, they connect interpretable early-warning alerts with a more complete assessment of conditional depeg and tail risk.


\subsubsection{Early-Warning Classification Model}
\label{sec:approach_classifier}

We formulate early warning of stablecoin depegs as an event-based binary classification problem. At each hour $t$, a classifier uses the available predictors $\mathbf{x}_t$ to estimate the probability of a new depeg event beginning between two and 24 hours ahead. The objective is to provide an actionable warning before the first threshold breach, rather than to mechanically identify an episode that is already underway.

\paragraph{Event definition and classification target}

Let $d_t$ denote the pool-price deviation from parity in basis points. A threshold breach occurs when $|d_t| \geq \delta$. The main analysis uses $\delta=15$ basis points for the USDC--USDT pool. Breaches are declustered into distinct events using a 24-hour reset period: a new event can begin only after the breach condition has been absent for 24 hours. Each event is represented by its first breach time, denoted by $\tau_e$. This prevents persistent deviations and closely spaced breaches from being treated as independent incidents.

The positive class comprises observations within the 24 hours preceding an event onset, subject to a minimum lead time of two hours. For retained observations, the target is
\begin{equation}
    y_t =
    \mathbb{I}\!\left\{
        \exists e:
        2\,\mathrm{h} \leq \tau_e-t \leq 24\,\mathrm{h}
    \right\}.
    \label{eq:early_warning_target}
\end{equation}

Already-depegged hours, for which $|d_t|\geq\delta$, and observations associated with threshold re-crossings that do not satisfy the event-reset criterion are censored. A new eligible event requires a preceding period of at least 24 consecutive hours below the depeg threshold; earlier re-crossings are treated as continuations of the existing episode. Censored observations are assigned neither positive nor negative labels: they do not represent opportunities to warn before a distinct event, but labeling them as negatives would incorrectly characterize ongoing depeg stress as a non-event condition. Hours inside the applicable minimum lead-time window are likewise censored. Positive labels therefore identify opportunities for advance warning of a new, eligible event.

The two-hour minimum lead time reduces the scope for learning a near-mechanical detection rule based on prices immediately approaching the breach threshold. It also accommodates the one-hour data-collection lag in our dataset: a two-hour separation between the predictor timestamp and event onset leaves at least one hour for an operational response. The maximum horizon of 24 hours allows time to investigate the warning and adjust liquidity or risk exposures.

\paragraph{Candidate models}

We compare four standard tree-based classifiers, XGBoost, LightGBM, Random Forest, and CatBoost,with two tabular foundation models, Kumo Tabular and TabICLv2 \cite{kumo2026tabular,qu2026tabiclv2}. The foundation models use pretrained representations and labeled historical observations as context for prediction. The inclusion of these specific models is motivated by their reported performance on tabular benchmarks, especially in non-iid settings. Their efficient inference relative to other tabular foundation models and publicly available model weights are also desirable for this task.

The small number of distinct depeg events provides a further motivation for the use of foundation models. Although an event generates several positive hourly labels, these observations do not represent independent examples of peg-loss risk. Pretraining may therefore provide useful inductive biases when relevant event examples are scarce. Whether these benefits extend to this highly imbalanced early-warning task is assessed empirically (section \ref{sec:cv_results}).

\paragraph{Chronological cross-validation and final holdout}

All observations from 1 January 2025 onward are reserved as a final holdout test set, matching the forecasting-model evaluation. This period is excluded from model selection and alert-threshold tuning.

Model comparison uses four chronological cross-validation folds with non-overlapping outer test windows covering 2024. The windows are balanced by event count rather than calendar duration, with at least four eligible depeg events in each fold.
Table~\ref{tab:classifier_cv} reports their boundaries and event counts. Together, the outer test windows contain 18 distinct eligible events, each evaluated in one outer test fold.

Within each fold, the fitting period precedes a validation period,
which in turn precedes the outer test period. A 24-hour embargo
separates fitting from validation and validation from testing,
preventing the forward-looking target from incorporating outcomes
in the next partition. Within each fold, the mean log-liquidity curve is estimated from the fitting partition only and held fixed throughout validation and testing. For foundation models, labeled context
observations are restricted to the corresponding fitting partition.
For each candidate specification, the operating threshold is
calibrated on the validation period and then fixed for evaluation
on the outer test period. Thus, fold-level evaluation includes both
threshold calibration and the subsequent alert-construction policy.

\begin{table}[pos=!htbp,align=\centering]
    \centering
    \small
    \caption{Outer test windows for the early-warning classifier.
    All timestamps are in UTC.}
    \label{tab:classifier_cv}
    \begin{tabular}{cllr}
        \hline
        Fold & Test start & Test end & Eligible events \\
        \hline
        1 & 2024-01-01 00:00 & 2024-03-06 08:00 & 4 \\
        2 & 2024-03-06 09:00 & 2024-07-01 02:00 & 4 \\
        3 & 2024-07-01 03:00 & 2024-11-21 18:00 & 5 \\
        4 & 2024-11-21 19:00 & 2024-12-31 23:00 & 5 \\
        \hline
        Total & & & 18 \\
        \hline
    \end{tabular}
\end{table}

\paragraph{Alert construction and event-level evaluation}

A probability threshold converts hourly predictions into an alerting
rule. Alert episodes use a fixed 12-hour quiet-period reset: after an
episode begins, subsequent threshold crossings belong to the same
episode until the score has remained below the threshold for
12 consecutive observed hours. Each episode is timestamped at its
initial crossing. The first episode in each evaluation partition
also requires this observed quiet history; missing timestamps or
scores interrupt the quiet period. False alert counts therefore measure distinct episodes rather than hourly threshold exceedances.

An event beginning at $\tau_e$ is detected in time if an alert episode
begins within $[\,\tau_e-24\,\mathrm{h},\ \tau_e-2\,\mathrm{h}\,]$.
Eligible events require complete observed prediction coverage over
this warning window. Each event contributes at most one successful
detection, irrespective of the number of timely episodes. Timely
event recall is thus
\begin{equation}
    R_{\mathrm{event}}
    =
    \frac{1}{|E|}
    \sum_{e\in E}
    \mathbb{1}\!\left\{
        \exists a_j:
        2\,\mathrm{h}\leq\tau_e-a_j\leq24\,\mathrm{h}
    \right\},
    \label{eq:timely_event_recall}
\end{equation}
where $E$ is the set of eligible events and $a_j$ denotes an
alert-episode onset.

False-alert burden is computed as
\begin{equation}
    F
    =
    730.5\,
    \frac{N_{\mathrm{false}}}{H_{\mathrm{observed}}},
    \label{eq:false_alert_burden}
\end{equation}
where $H_{\mathrm{observed}}$ counts score hours with complete future
follow-up and 730.5 is the number of hours in an average calendar
month. We also report the median lead time of the first timely alert among detected events.

\paragraph{Neyman--Pearson threshold calibration}

Operating thresholds are calibrated using the Neyman--Pearson (NP) \citep{neyman1933}
order-statistic rule. Let $\eta$ denote the target conditional hourly false-positive rate (FPR). For each candidate model and fold, calibration uses validation scores from eligible negative hours. Retained negative observations are separated by at least 24 hours in order to reduce the effect of serial dependence.

Let $s_{(1)}\leq\cdots\leq s_{(n)}$ denote the retained negative scores.
For a fixed nominal violation tolerance $\delta_{\mathrm{NP}}=0.05$,
the selected rank is
\begin{equation}
    k^\star
    =
    \min\left\{
        k\in\{1,\ldots,n\}:
        \sum_{j=k}^{n}
        \binom{n}{j}(1-\eta)^j\eta^{n-j}
        \leq\delta_{\mathrm{NP}}
    \right\}.
    \label{eq:np_rank_selection}
\end{equation}

The operating rule alerts on scores strictly exceeding $s_{(k^\star)}$. If no admissible rank exists, the fold uses an infinite cutoff and emits no alerts. The cutoff is fixed before evaluation on the outer test period.

Under iid negative calibration scores independent of the fitted scoring function, this construction controls the probability that the population hourly false-positive rate exceeds $\eta$.
These assumptions are not established for the present time series: residual serial dependence may remain after spacing, and validation data used for early stopping are not independent calibration data. Accordingly, the NP bound is treated as nominal, and the resulting policy is assessed through chronological out-of-sample evaluation.

The hourly FPR target is not equivalent to a budget of false-alert episodes per month. Their relationship depends on the temporal pattern of threshold exceedances and the episode-reset rule. The choice of $\eta$ is discussed in section \ref{sec:cv_results}, with sensitivity to alternative targets reported in the appendix \ref{sec:policy_sensitivity}.

\paragraph{Operational model selection and uncertainty}

The operational constraint is five false-alert episodes per month. For each candidate specification, event recall is pooled over the 18 distinct outer-test events, while false-alert burden pools false episode counts and observed follow-up hours across folds. Recall and median-lead confidence intervals resample unique realized events;
false-alert-burden intervals resample calendar blocks.

Let $U_F(m,\alpha;\eta)$ denote the upper endpoint of the bootstrap
95\% confidence interval for pooled false-alert burden. At a given
NP target, model selection follows
\begin{equation}
    (m^\star,\alpha^\star)
    \in
    \underset{
        (m,\alpha):
        U_F(m,\alpha;\eta)\leq5
    }{\operatorname{arg\,max}}
    \ R_{\mathrm{event,CV}}(m,\alpha;\eta).
    \label{eq:operational_model_selection}
\end{equation}
Lower pooled false-alert burden breaks ties in event recall.
If no candidate satisfies the constraint, no specification is declared operationally admissible at that target.

This criterion prioritizes detecting distinct events with usable advance notice while accounting for uncertainty in alert burden. The bootstrap upper endpoint is an operational selection criterion, not a guarantee that future average burden, or the burden in every individual month, remains below five episodes.

ROC AUC, precision--recall AUC and Brier score are reported as secondary measures of discrimination and probability accuracy, rather than selection criteria. Comparisons across model families, $\alpha$ values and NP targets constitute development-stage model selection. The selected specification is subsequently evaluated on the reserved holdout. Results are presented in Section~\ref{sec:cv_results}.


\subsubsection{Probabilistic Forecasting Model}
While the early-warning classifier provides a binary alert, it does not convey the magnitude or distributional shape of potential price deviations. Risk managers and DeFi participants often need to reason about tail risk, for example, the probability that a deviation will exceed a collateral liquidation threshold, or the expected shortfall conditional on a depeg occurring. To address this, we complement the classifier with a probabilistic forecasting model that outputs the full conditional distribution $p(P_{[t,\,t+T]} \mid \mathbf{x}_t)$
of the pool price over the same 24-hour horizon. This distribution-level output allows users to assess not only whether a depeg is likely but how severe it might be, providing a more granular and practically useful characterization of tail risk than a binary signal alone. The two models are designed to be used together: the classifier provides an interpretable alert, while the probabilistic model provides the distributional context needed to size exposures and set risk limits accordingly.

\begin{figure}[pos=!htbp,align=\centering]
\centering
\includegraphics[width=\linewidth]{figures/archi.png}
\caption{Overview of the neural network architecture.\href{https://github.com/QuantLet/onchain-insights/tree/main/11.\%20Forecasting\%20architecture}{\includegraphics[height=0.6\baselineskip]{figures/Quantlet LOGO.png}}}
\label{fig:diagram}
\end{figure}

\paragraph{Neural Architecture Design}
We design a neural architecture for univariate probabilistic forecasting with historical exogenous features, balancing predictive performance, computational efficiency, and interpretability (Figure~\ref{fig:diagram}). Following recent related forecasting architectures, separate branches encode the target series (depeg deviation in basis points) and exogenous covariates \citep{TIDE,TimeXer}. 

Both branches apply Reversible Instance Normalization (ReVIN) \citep{revin} to their inputs. After forecasting, the nonstationary information in the price series is reinjected through a learnable affine transformation. 

The target branch employs dilated causal temporal convolution to capture dynamics across multiple time scales without applying attention over the full lookback window \citep{oord2016wavenet,bai2018tcn,chen2020deeptcn,moreno2024deepvol,chen2024financialcrisis_tcn}. Its output is mapped to one latent embedding vector per forecast horizon, limiting subsequent attention computation to the shorter-horizon representation.

The covariate branch maps each exogenous trajectory to a variable-level embedding vector \citep{itransformer}. A context-conditioned selection module scores covariates with a Gated Residual Network (GRN) \citep{TFT}. Cross-attention layers then enable target horizon tokens to attend to the selected covariate tokens, and an output head maps the fused representations to distributional parameters. 

The selection module is sparse by design, producing a weight for each exogenous feature in each input window, with sparsemax activations assigning exact zeros to features excluded from the selected representation. We further add the mean Shannon entropy of the feature-selection weights to the loss, encouraging the model to concentrate weight on fewer exogenous features. The selection weights and horizon-wise attention maps (Figure~\ref{fig:variate_selection}) provide diagnostics of the exogenous inputs emphasized by the model, complementing assessments of predictive accuracy and calibration. Hyperparameters for the reported runs are provided in Appendix~\ref{sec:hyperparameters}.

\paragraph{Distributional Forecast Head} 
Concentrated liquidity in stablecoin AMMs can produce small fluctuations around the peg alongside rare, large deviations, as examined in the residual analysis in Appendix~\ref{sec:residual_analysis}. To accommodate heavy tails and potential asymmetry between USDC and USDT depegs, the forecast head parameterizes a quantile function with a flexible central region and separate generalized Pareto (GPD) tails. The central quantile function, $Q_t^{B}$, uses an I-spline basis with nonnegative, softplus-constrained coefficients to maintain monotonicity \citep{xu2023}. For each forecast horizon, the network outputs a body-median anchor $\ell_t$, a positive scale $s_t$, spline coefficients $\tilde a_{t,1:K}$, and lower- and upper-tail shape parameters $\xi_L$ and $\xi_U$. Specifically, let:
\begin{equation}
Q_t^0(u)=\sum_{k=1}^{K} a_{t,k} I_k(u),
\qquad
a_{t,k}=\operatorname{softplus}(\tilde a_{t,k})+\epsilon,
\end{equation}
where $I_k$ denotes the $k$th I-spline basis function. The body shape is
anchored at its median,
\begin{equation}
C_t(u)=Q_t^0(u)-Q_t^0(0.5).
\end{equation}
Its variance $V_t$ is estimated on the quantile grid $\{u_j\}_{j=1}^{N}$
using quadrature weights $\{\omega_j\}_{j=1}^{N}$:
\begin{equation}
\bar C_t=\sum_{j=1}^{N}\omega_j C_t(u_j),
\qquad
V_t=\sum_{j=1}^{N}\omega_j
       \bigl(C_t(u_j)-\bar C_t\bigr)^2.
\end{equation}
The calibrated body is:
\begin{equation}
Q_t^B(u)=\ell_t+
          \frac{s_t}{\sqrt{V_t+\epsilon}}\,C_t(u).
\end{equation}
Thus, $\ell_t$ anchors the body median and $s_t$ sets its grid-based
standard deviation, up to numerical regularization.

The lower and upper GPD tails are joined to the body at the splice
quantiles $\tau_L$ and $\tau_U$:
\begin{equation}
Q_t(u)=
\begin{cases}
z_L-\dfrac{\beta_L}{\xi_L}
\left[\left(\dfrac{u}{\tau_L}\right)^{-\xi_L}-1\right],
& 0<u<\tau_L,\\[2ex]
Q_t^B(u),
& \tau_L\leq u\leq\tau_U,\\[1.2ex]
z_U+\dfrac{\beta_U}{\xi_U}
\left[\left(\dfrac{1-u}{1-\tau_U}\right)^{-\xi_U}-1\right],
& \tau_U<u<1,
\end{cases}
\end{equation}
where
\[
z_L=Q_t^B(\tau_L),\qquad z_U=Q_t^B(\tau_U),\qquad
\beta_L=\tau_L(Q_t^B)'(\tau_L),\qquad
\beta_U=(1-\tau_U)(Q_t^B)'(\tau_U).
\]
The tail scales therefore match the body derivatives at the splice
points. The shape parameters are constrained to the ranges reported in
Table~\ref{tab:angel_hyperparameters}; when a shape parameter is zero,
the corresponding expression is interpreted by its continuous
exponential limit.

The location and scale calibration applies to the body before the GPD tails are spliced. Consequently, $\ell_t$ and $s_t$ are not in general the mean and standard deviation of the full spliced distribution; for $\xi\geq 1/2$, a GPD tail has infinite variance. The spliced parameterization supports evaluation of tail-risk relevant quantities such as expected shortfall.

\paragraph{Training objective}

We train the predictive distribution using the continuous ranked probability score (CRPS) \citep{CRPS} in its well known approximation on a finite quantile grid $\mathcal{U}=\{u_m\}_{m=1}^{M}$:

\begin{equation} \widehat{\operatorname{CRPS}}(F_t,y)=
2\sum_{m=1}^{M} \eta_m\, \rho_{u_m}\left(y-Q_t(u_m)\right), \end{equation}

where $\eta_m$ are quadrature weights. We additionally use threshold-weighted CRPS to emphasize price regions associated with depegs \citep{twCRPS}. Its quantile-grid approximation is 

\begin{equation} \widehat{\operatorname{twCRPS}}(F_t,z; w)=
2\sum_{m=1}^{M}
\eta_m\,
\rho_{u_m}
\left(
v(z)-v(Q_t(u_m))
\right)\,
\qquad
v(r)=\int_{0}^{r} w(s)\,ds,
\end{equation}
where $w(r)\geq 0$ is a weight function that emphasizes regions of the outcome space. In our setting, we choose a smooth two-sided threshold weight that focuses the loss on deviations with significant absolute magnitude:

\begin{equation}
w(r)
=
\sigma\!\left(\frac{r-\delta}{h}\right)
+
\sigma\!\left(\frac{-r-\delta}{h}\right),
\end{equation}



\subsection{Economic evaluation: a liquidity-provision backtest}\label{sec:lp_method}

We evaluate both models in the decision problem most exposed to the pool's depeg risk: providing liquidity
to it. A liquidity provider (LP) runs an automated vault in the USDC--USDT 0.01\% pool, like the Gamma vaults
that manage Uniswap V3 positions for their depositors. Every hour the vault decides the price range in which
its capital is deployed. A narrow range earns a larger share of swap fees, but the vault must reposition it
more often, paying gas and swap costs, and every price move inside the range is exploited by arbitrageurs
\citep{Milionis2022}. In a depeg, a vault that keeps recentring on the moving price repeatedly sells the
appreciating token and buys the depreciating one.

\paragraph{Position value, fees and loss-versus-rebalancing.} Let $P_k$ be the pool price of USDC in USDT at
the end of hour $k$. A position with liquidity $L$ on $[p_a,p_b]$ holds
$x(P)=L\big(1/\sqrt{\bar P}-1/\sqrt{p_b}\big)$ USDC and $y(P)=L\big(\sqrt{\bar P}-\sqrt{p_a}\big)$ USDT, with
$\bar P=\min\{\max\{P,p_a\},p_b\}$ and value $V(P)=x(P)P+y(P)$ \citep{uniswap_whitepaper}. Its fee income in
hour $k$ is
\begin{equation}
f_k=F_k\,\omega_k\,\frac{L}{L^{\mathrm{act}}_k+L},
\end{equation}
where $F_k$ is the pool's fee revenue, $L^{\mathrm{act}}_k$ its active liquidity and $\omega_k$ the share of
the hour's high--low price interval inside the range. The benchmark is the \emph{rebalancing portfolio} of
\citet{Milionis2022}: it holds the same USDC inventory as the position at the start of each hour and trades
at market prices. The position's loss relative to it,
\begin{equation}
\mathrm{LVR}_k=V(P_k)-V(P_{k-1})-x(P_{k-1})\,(P_k-P_{k-1})\le 0,
\end{equation}
is the price the LP pays to arbitrageurs, and the vault's net profit relative to the rebalancing portfolio is
\begin{equation}
\Pi_k=f_k+\mathrm{LVR}_k-c_k,
\end{equation}
where $c_k$ is the cost of repositioning: \$10 of gas per withdrawal and deposit, and the pool fee plus
0.5\,bp of slippage on the swap that restores the inventory required by the new range. The vault's value
satisfies $V_T-V_0=\sum_k\Pi_k+\sum_k X_{k-1}(P_k-P_{k-1})$, where $X$ is its total USDC holding; the second
term is the exposure to the USDC/USDT price, which the rebalancing portfolio shares, so $\Pi$ isolates the
consequences of the LP's decisions. All results are for a \$1M vault, in basis points of its capital.

\paragraph{Vaults.} Two standard vaults serve as benchmarks: a static range of $\pm5$ ticks (about
$\pm5$\,bp) around parity, never repositioned, and a Gamma-style vault that recentres a $\pm2$-tick range on
the current price whenever the price leaves it. The $\pm2$-tick width is the best recentring width both
before and during the holdout. Each model then has a single role, with no parameter tuned on the holdout.
\begin{itemize}
  \item \emph{Forecaster: range sizing.} The range is the forecaster's 10\%--90\% predictive interval, taken
        as the union of the intervals for the next six hours and extended to contain the current price, so it
        widens when uncertainty rises and leans towards the side where the forecast distribution places more
        mass. The vault repositions when the price leaves the range or when the target width becomes more than
        twice, or less than a third of, the current width.
  \item \emph{Classifier: pause during alerts.} During an alert episode, defined by the operating cutoff
        and 12-hour quiet reset of Section~\ref{sec:holdout_results}, the vault keeps its position and does
        not reposition. A position left in place leaves the range as the price moves away and then stops
        incurring LVR, instead of being recentred on the moving price. The classifier is used at the 15\,bp
        threshold for which it is trained.
\end{itemize}
The two-model vault uses the forecaster's ranges and pauses during the classifier's alerts. The same pause
is applied to the Gamma-style vault to isolate the classifier's contribution.

\paragraph{Controls.} Three controls separate the information in the classifier from the act of pausing. A
\emph{price trigger} pauses the vault while $|d_k|\ge15$\,bp, with the same quiet reset; it acts on the
events the classifier is trained to anticipate, but only once they have started. An \emph{oracle} pauses the
vault from 24 hours before each event. \emph{Placebo} schedules place the classifier's \lpEpisodes{} alert
episodes, with their durations, at \lpNPlacebo{} random sets of dates. Decisions are taken at the end of hour
$k$ for hour $k+1$; price rules use $P_k$, while model outputs are used with a one-hour data lag.

\paragraph{Performance and risk measures.} We report fee income, LVR, costs and net profit (annualised
$\sum_k\Pi_k$), the Sharpe ratio of daily and the Sortino ratio of hourly $\Pi$, the maximum drawdown of
cumulative $\Pi$, the worst day, and the 5\% conditional value-at-risk (CVaR) of daily $\Pi$. Because the
classifier targets rare events, we also measure risk where it materialises: the \emph{event drawdown} is
the sum, over the 15\,bp events, of the largest drawdown of cumulative $\Pi$ between 24 hours before and
48 hours after each event.

\paragraph{Samples.} The main sample is the common holdout coverage of both models, from \lpHoldoutStart{}
to \lpHoldoutEnd, with \lpNEvents{} events. As an out-of-sample check we use the classifier's 2024
cross-validation folds. The pool's subgraph records no fees before 19 March 2024, so this sample starts on
20 March 2024 and contains \lpYNEvents{} events; the neural forecasts cover only the holdout, so a rolling
empirical forecaster (the trailing 30-day distribution of $h$-hour price changes) stands in for them.

\section{Results}\label{sec:results}

\subsection{Early-Warning Classification Model}
\subsubsection{Cross-validated operational performance} \label{sec:cv_results}

Table~\ref{tab:cv_model_alpha_results} compares six classifier families across four Gegenbauer $\alpha$ values at the 15-basis-point depeg threshold. Operating cutoffs are calibrated separately within each fold using the NP rule with hourly FPR target $\eta=0.08$ and nominal violation tolerance $\delta_{\mathrm{NP}}=0.05$. The operating target is selected through development-stage cross-validation; sensitivity to alternative threshold-selection policies is reported in Appendix~\ref{sec:policy_sensitivity}. The hourly FPR target is distinct from the operational constraint of five false-alert episodes per month.

The foundation model Kumo with $\alpha=1$ achieves the highest pooled timely event recall, detecting 11 of the 18 eligible events ($61.1\%$). Its false-alert burden is 3.11 episodes per month (95\% CI: 1.90--4.46). Because the upper confidence endpoint remains below five episodes per month, this specification satisfies the operational selection criterion and is selected for subsequent holdout evaluation. Its median first-alert lead time is 4 hours (95\% CI: 3--13 hours), indicating that successful warnings typically occur toward the later part of the admissible warning window.

Tree-based classifiers heavily underperform the foundation models on this task, both in row-level ranking metrics, particularly ROC AUC, and in operational event detection. The strongest tree-based specification, CatBoost with $\alpha=0.1$, detects 5/18 events, fewer than half the 11/18 detected by the selected Kumo model. Several tree-based specifications detect no eligible events, despite generating false alerts.

The sensitivity analysis in Appendix~\ref{sec:policy_sensitivity} also selects a foundation model under every tested operating policy, suggesting that their advantage is not confined to a single threshold-calibration setting. These findings motivate the use of foundation models in this application, where the scarcity of realized depeg events limits the information available for task-specific model fitting.

\begin{table}[pos=!htbp,align=\centering]
\caption{Operational out-of-sample comparison at the 15-basis-point depeg threshold. Probability cutoffs are selected on preceding validation data using the Neyman--Pearson order-statistic rule with hourly FPR limit $\eta=0.08$ and nominal violation tolerance $\delta_{\mathrm{NP}}=0.05$. Models are selected by event-level recall, under an operational constraint of 5 false alerts per month. Timely warnings precede onset by 2–24 hours. Recall and pooled median-lead confidence intervals resample unique realized events; alert-burden intervals resample calendar blocks. Best displayed values in each metric are highlighted in bold red and second-best distinct displayed values in underlined blue.}
\label{tab:cv_model_alpha_results}
\resizebox{\textwidth}{!}{
\begin{tabular}{llrrrrrr}
\toprule
Model & $\alpha$ & Event recall & False alerts/month & Median lead (h) & AUC & AUPRC & Brier score \\
\midrule
\multirow{4}{*}{Catboost} & 0.1 & 0.278 (5/18) & 0.92 [0.33, 1.57] & 13.0 [8.0, 19.0] & 0.643 & 0.199 & 0.172 \\
 & 0.3 & 0.167 (3/18) & 0.67 [0.09, 1.32] & 13.0 [13.0, 17.0] & 0.594 & 0.103 & 0.165 \\
 & 0.5 & 0.222 (4/18) & 1.01 [0.33, 1.83] & 12.0 [6.0, 18.0] & 0.610 & 0.137 & 0.148 \\
 & 1 & 0.167 (3/18) & 0.76 [0.25, 1.33] & 12.0 [2.0, 12.0] & 0.584 & 0.119 & 0.167 \\
\midrule
\multirow{4}{*}{XGBoost} & 0.1 & 0.167 (3/18) & 1.18 [0.50, 1.95] & 18.0 [4.0, 19.0] & 0.583 & 0.130 & 0.239 \\
 & 0.3 & 0.056 (1/18) & 1.09 [0.42, 1.83] & \textcolor{blue}{\underline{19.0 [19.0, 19.0]}} & 0.621 & 0.131 & 0.241 \\
 & 0.5 & 0.167 (3/18) & 1.09 [0.42, 1.83] & 13.0 [12.0, 18.0] & 0.637 & 0.141 & 0.227 \\
 & 1 & 0.056 (1/18) & 0.50 [0.00, 1.23] & \textcolor{blue}{\underline{19.0 [19.0, 19.0]}} & 0.636 & 0.119 & 0.194 \\
\midrule
\multirow{4}{*}{LightGBM} & 0.1 & 0.000 (0/18) & \textcolor{red}{\textbf{0.25 [0.00, 0.69]}} & NA & 0.598 & 0.114 & \textcolor{blue}{\underline{0.062}} \\
 & 0.3 & 0.000 (0/18) & \textcolor{red}{\textbf{0.25 [0.00, 0.69]}} & NA & 0.530 & 0.109 & \textcolor{blue}{\underline{0.062}} \\
 & 0.5 & 0.000 (0/18) & \textcolor{red}{\textbf{0.25 [0.00, 0.69]}} & NA & 0.525 & 0.116 & \textcolor{blue}{\underline{0.062}} \\
 & 1 & 0.000 (0/18) & \textcolor{blue}{\underline{0.34 [0.00, 0.77]}} & NA & 0.520 & 0.129 & \textcolor{blue}{\underline{0.062}} \\
\midrule
\multirow{4}{*}{Random Forest} & 0.1 & 0.000 (0/18) & \textcolor{blue}{\underline{0.34 [0.00, 0.76]}} & NA & 0.677 & 0.180 & \textcolor{red}{\textbf{0.061}} \\
 & 0.3 & 0.111 (2/18) & 1.01 [0.25, 1.84] & \textcolor{red}{\textbf{19.5 [19.0, 20.0]}} & 0.669 & 0.189 & \textcolor{red}{\textbf{0.061}} \\
 & 0.5 & 0.056 (1/18) & 0.59 [0.17, 1.10] & \textcolor{blue}{\underline{19.0 [19.0, 19.0]}} & 0.671 & 0.160 & \textcolor{red}{\textbf{0.061}} \\
 & 1 & 0.111 (2/18) & \textcolor{red}{\textbf{0.25 [0.00, 0.51]}} & 12.5 [12.0, 13.0] & 0.637 & 0.108 & 0.064 \\
\midrule
\multirow{4}{*}{Kumo} & 0.1 & 0.444 (8/18) & 3.87 [2.59, 5.23] & 4.0 [3.0, 12.0] & 0.770 & 0.208 & 0.068 \\
 & 0.3 & 0.389 (7/18) & 3.70 [2.43, 5.04] & 5.0 [2.0, 21.0] & 0.781 & 0.202 & 0.067 \\
 & 0.5 & \textcolor{blue}{\underline{0.556 (10/18)}} & 3.95 [2.74, 5.22] & 9.5 [3.0, 19.0] & 0.778 & 0.204 & 0.068 \\
 & 1 & \textcolor{red}{\textbf{0.611 (11/18)}} & 3.11 [1.90, 4.46] & 4.0 [3.0, 13.0] & \textcolor{red}{\textbf{0.797}} & 0.204 & 0.064 \\
\midrule
\multirow{4}{*}{TabICL-v2} & 0.1 & 0.444 (8/18) & 2.86 [1.67, 4.19] & 8.5 [6.0, 21.0] & \textcolor{blue}{\underline{0.790}} & \textcolor{blue}{\underline{0.227}} & 0.079 \\
 & 0.3 & 0.389 (7/18) & 3.11 [1.91, 4.39] & 7.0 [4.0, 21.0] & 0.775 & 0.211 & 0.076 \\
 & 0.5 & 0.389 (7/18) & 2.19 [1.34, 3.07] & 7.0 [5.0, 21.0] & 0.750 & 0.203 & 0.083 \\
 & 1 & 0.167 (3/18) & 2.61 [1.65, 3.68] & 12.0 [4.0, 21.0] & 0.778 & \textcolor{red}{\textbf{0.252}} & 0.076 \\
\bottomrule
\end{tabular}
}
\end{table}

\subsubsection{Final holdout performance}
\label{sec:holdout_results}

The CV-selected Kumo specification with Gegenbauer $\alpha=1$ is evaluated on the reserved holdout test set from 1 January 2025 to 6 May 2026. The NP operating cutoff is fixed before holdout evaluation.

The model detects all five eligible depeg events. Their first timely alerts precede onset by 3, 22, 2, 9 and 23 hours. 

Although all five events are detected, the small denominator limits inference about the detection rate in future periods. The reported recall should thus not be interpreted as strong evidence of certain detection.

The model generates 39 false-alert episodes over 11,741 hours, corresponding to 2.43 episodes per month (calendar-block bootstrap 95\% CI: 1.56--3.37). Both the point estimate and upper confidence endpoint are well below the operational budget of five episodes per month. False alerts remain temporally concentrated, including periods during March--April and September--November 2025 (Figure \ref{fig:holdout_policy}).

Row-level discrimination is also stronger than in the CV comparison: holdout ROC AUC is 0.910 and AUPRC is 0.103, compared with CV fold means of 0.797 and 0.204, respectively. The differing directions of these metrics emphasize that their values depend on the evaluation population and class prevalence as well as score ranking. The holdout Brier score is 0.012.

The observed hourly FPR is 4.97\%, below the nominal 8\% calibration target. Hourly TPR is 30.4\%, despite detection of every eligible event. These quantities describe different objectives: event recall requires at least one timely episode per event, whereas hourly TPR measures the fraction of positive warning-window hours classified above the cutoff. Successful event detection therefore does not require continuous warning throughout the lead window. Figure~\ref{fig:holdout_policy} illustrates the timing of alerts relative to the realized event onsets.

\begin{figure}[pos=!htbp,align=\centering]
    \centering
    \includegraphics[width=\textwidth]{figures/holdout_policy.png}
    \caption{Holdout alert outcomes for the CV-selected Kumo specification with $\alpha=1$ and NP hourly FPR target $\eta=0.08$ at the 15-basis-point depeg threshold. Issued alerts are highlighted, with red crosses indicating a false alert and green triangles a timely depeg detection.\href{https://github.com/QuantLet/onchain-insights/tree/main/8.\%20Early-Warning\%20Classifier\%20Cross\%20Validation}{\includegraphics[height=0.6\baselineskip]{figures/Quantlet LOGO.png}}}
    \label{fig:holdout_policy}
\end{figure}

\subsubsection{Economic feature-block ablations}
\label{sec:ablation_results}

\begin{table}[pos=!htbp,align=\centering]
\caption{Cross-validated feature-block ablations. All variants use the full model's hyperparameters and the same operating-threshold selection procedure.}
\label{tab:ablation_15_kumo_tabular}
\begin{tabular}{llllrrr}
\toprule
Ablated Set & Event recall & False alerts/month & Median lead (h) & AUC & AUPRC \\
\midrule
None & 0.611 (11/18) & 3.11 [1.90, 4.46] & 4.0 [3.0, 13.0] & 0.797 & 0.204 \\
All but price & 0.722 (13/18) & 7.23 [5.83, 8.76] & 8.0 [5.0, 17.0] & 0.727 & 0.182 \\
Curve 3pool liquidity conditions & 0.333 (6/18) & 4.12 [2.89, 5.47] & 7.5 [3.5, 21.0] & 0.725 & 0.140  \\
Liquidity ownership, position structure & 0.444 (8/18) & 2.52 [1.51, 3.62] & 6.5 [4.0, 21.0] & 0.770 & 0.152  \\
Broader market conditions & 0.444 (8/18) & 3.53 [2.32, 4.68] & 7.5 [3.0, 13.0] & 0.785 & 0.219\\
Historical peg deviation & 0.500 (9/18) & 4.29 [3.08, 5.49] & 8.0 [3.0, 14.0] & 0.746 & 0.149 \\
Trading velocity and flows & 0.389 (7/18) & 3.87 [2.70, 5.20] & 9.0 [3.0, 21.0] & 0.768 & 0.156 \\
Lending market conditions & 0.333 (6/18) & 4.37 [3.00, 5.85] & 7.0 [3.5, 19.0] & 0.737 & 0.162 \\
Uniswap liquidity curve Shape & 0.222 (4/18) & 4.04 [2.81, 5.46] & 6.5 [3.0, 15.0] & 0.682 & 0.148 \\
\bottomrule
\end{tabular}
\end{table}

We examine the contribution of feature blocks corresponding to the depeg channels identified in Section~\ref{sec:background}. The full feature list and group assignments are provided in Appendix~\ref{sec:data-appendix}.

Table~\ref{tab:ablation_15_kumo_tabular} supports the use of the combined feature set. The full specification achieves the second-highest event recall, detecting 11/18 events, alongside the highest ROC AUC and second-highest AUPRC. The price-only reference detects 13/18 events but generates 7.23 false alerts per month, with even its lower confidence bound exceeding the five-episode operating limit. Its row-level ranking metrics also deteriorate relative to the full specification.

Only the liquidity-ownership and broader-market ablations remain admissible under the false-alert upper-confidence-bound criterion. These variants achieve the third- and second-highest ROC AUC, respectively, but each detects only 8/18 events.

Removing Uniswap liquidity-curve features produces the largest deterioration, reducing detection to 4/18 events and ROC AUC to 0.682. Removing Curve 3pool liquidity conditions reduces detection to 6/18 events. Both variants have higher false-alert burden than the full model and lower AUC and AUPRC than the price-only reference. Removing AAVE lending conditions likewise reduces detection to 6/18 events, indicating a substantial complementary contribution.

Overall, the results support the predictive relevance of the identified depeg channels. Uniswap V3 liquidity-curve shape and Curve 3pool liquidity conditions provide particularly informative signals, while trading flows, lending conditions and broader market conditions contribute complementary risk information.

\subsection{Probabilistic Forecasting Model}
\subsubsection{Forecasting setup and evaluation design}

The probabilistic forecasting model is evaluated through full distributional forecasts of the USDC/USDT peg deviation over a 24-hour horizon. The holdout test period matches that of the early-warning model and spans January 2025 to May 2026. Forecast quality is assessed using the continuous ranked probability score (CRPS) and the threshold-weighted CRPS (twCRPS), which also serve as training objectives for the neural models. These metrics capture, respectively, overall distributional accuracy and performance in regions most relevant to depeg risk. We additionally report quantile losses at the 1\%, 5\%, 95\%, and 99\% quantiles, and compute calibration diagnostics to assess the reliability of the predictive distributions. To assess whether differences in aggregate forecast scores are statistically significant, we apply one-sided Diebold--Mariano tests to paired CRPS and twCRPS losses; 95\% confidence intervals for mean loss differences are computed using a block bootstrap to account for serial dependence in the loss differentials.

To contextualize performance, we compare the proposed model against several benchmarks: a naive random-walk model with empirical residuals, a standard ARIMA specification, and a GARCH model with Student-$t$ innovations. We also include two neural forecasting architectures suited to univariate forecasting with historical exogenous regressors, namely TiDE \citep{TIDE} and the transformer-based TimeXer \citep{TimeXer}. To isolate the contribution of the forecasting backbone, both neural benchmarks are equipped with the same probabilistic head and training objective as the proposed model, based on a spline quantile function with spliced GPD tails and the twCRPS loss.

\subsubsection{Aggregate distributional forecasting performance}

Table~\ref{tab:benchmark} reports the aggregate out-of-sample distributional forecasting performance over the 24-hour forecast horizon. The proposed model achieves the lowest loss on every reported metric. Relative to the strongest benchmark for each metric, it improves CRPS by $6.6\%$ over TimeXer and twCRPS by $14.2\%$ over TiDE. The twCRPS improvement is more than twice the CRPS improvement, indicating that our model's gains are particularly strong in regions relevant to depeg risk. The lower- and upper-tail twCRPS improvements over TiDE are $8.9\%$ and $15.3\%$, respectively.

A similar pattern appears in the quantile losses. Relative to TimeXer, our model improves QL at the 1\% and 99\% quantiles by approximately $13.0\%$ and $10.8\%$, compared with gains of $6.5\%$ at the 5\% quantile and $5.3\%$ at the 95\% quantile. Thus, the improvements at the most extreme quantiles are about twice those at the 5\% and 95\% quantiles. These results suggest that our model improves overall distributional accuracy while providing especially strong gains for rare events and large deviations from parity.

Table~\ref{tab:model_benchmark_comparison} assesses the statistical significance of the CRPS and twCRPS differences. Mean loss differences are negative for all five benchmarks on both metrics, and every 95\% block-bootstrap confidence interval lies below zero. All one-sided Diebold--Mariano $p$-values are below $0.05$, indicating statistically significant outperformance by our model on both aggregate scores against every benchmark.

The classical benchmarks perform substantially worse in aggregate. GARCH improves on ARIMA and the naive forecast, as expected given its ability to model conditional heteroskedasticity, but it remains behind the neural models. This is consistent with the residual diagnostics motivating a flexible distributional approach: volatility modeling alone is insufficient to capture the asymmetric and state-dependent tail behavior of stablecoin peg deviations. Overall, the results support the use of probabilistic neural forecasting for this task and show that our models's strongest gains occur in the distributional regions most relevant to depeg-risk monitoring.

\begin{table}[pos=!htbp,align=\centering]
\centering
\caption{Out-of-sample distributional forecasting performance aggregated over the 24-hour forecast horizon. Lower values indicate better performance. CRPS measures aggregate distributional accuracy, while twCRPS emphasizes depeg-relevant regions of the outcome distribution. The low/high twCRPS components and extreme quantile losses evaluate lower- and upper-tail forecast quality separately.}
\label{tab:benchmark}
\resizebox{\textwidth}{!}{  
\begin{tabular}{lrrrrrrrr}
\toprule
Model & CRPS & twCRPS & $\operatorname{twCRPS}_{lower}$ & $\operatorname{twCRPS}_{higher}$ & QL 1\% & QL 5\% & QL 95\% & QL 99\% \\
\midrule
Our Model & \textcolor{red}{\textbf{0.8205}} & \textcolor{red}{\textbf{0.0665}} & \textcolor{red}{\textbf{0.0121}} & \textcolor{red}{\textbf{0.0545}} & \textcolor{red}{\textbf{0.0706}} & \textcolor{red}{\textbf{0.2081}} & \textcolor{red}{\textbf{0.1952}} & \textcolor{red}{\textbf{0.0611}} \\
GARCH & 2.2140 & 0.1060 & 0.0306 & 0.0754 & 0.1123 & 0.3884 & 0.3229 & 0.1057 \\
TimeXer & \textcolor{blue}{\underline{0.8787}} & 0.0883 & 0.0206 & 0.0677 & \textcolor{blue}{\underline{0.0811}} & \textcolor{blue}{\underline{0.2225}} & \textcolor{blue}{\underline{0.2062}} & \textcolor{blue}{\underline{0.0685}} \\
TiDE & 0.9702 & \textcolor{blue}{\underline{0.0775}} & \textcolor{blue}{\underline{0.0132}} & \textcolor{blue}{\underline{0.0643}} & 0.1560 & 0.3061 & 0.3009 & 0.1555 \\
ARIMA & 4.4917 & 1.3344 & 0.6747 & 0.6597 & 0.4212 & 1.4877 & 1.4975 & 0.4232 \\
Naive & 5.4997 & 2.1775 & 1.1169 & 1.0607 & 0.5190 & 1.8347 & 1.8347 & 0.5190 \\
\bottomrule
\end{tabular}
}
\end{table}


\begin{table}[pos=!htbp,align=\centering]
\centering
\caption{Comparison of our model performance against benchmark models aggregated over 24-hour
forecast horizon. Mean difference is defined as
$\mathrm{Loss}_{\text{our model}} - \mathrm{Loss}_{\text{benchmark}}$;
negative values favor our model. Improvement is the percentage reduction in
mean loss relative to the benchmark. Confidence intervals are bootstrap
confidence intervals. DM $p$-values are from one-sided tests of whether our
model has lower loss than the benchmark. Underlined $p$-values indicate that
our model significantly outperforms the corresponding benchmark at the 5\%
significance level.}
\label{tab:model_benchmark_comparison}
\small
\begin{tabular}{llrrr}
\toprule
Metric & Benchmark & Mean diff. (95\% CI) & Improvement (\%) & DM $p$-value \\
\midrule
\multirow{5}{*}{\textit{CRPS}}
 & GARCH & -1.3935 [-1.5789, -1.2159] & 62.9410 & $\underline{5.83 \times 10^{-89}}$ \\
 & TimeXer & -0.0582 [-0.0762, -0.0411] & 6.6231 & $\underline{5.36 \times 10^{-12}}$ \\
 & TiDE & -0.1497 [-0.1704, -0.1298] & 15.4321 & $\underline{7.98 \times 10^{-58}}$ \\
 & ARIMA & -3.6712 [-3.7344, -3.6010] & 81.7335 & $\underline{< 10^{-100}}$ \\
 & Naive & -4.6792 [-4.7461, -4.6077] & 85.0815 & $\underline{< 10^{-100}}$ \\
\addlinespace \midrule
\multirow{5}{*}{\textit{twCRPS}}
 & GARCH & -0.0395 [-0.0631, -0.0197] & 37.2463 & $\underline{1.34 \times 10^{-5}}$ \\
 & TimeXer & -0.0218 [-0.0332, -0.0124] & 24.6483 & $\underline{2.26 \times 10^{-6}}$ \\
 & TiDE & -0.0110 [-0.0181, -0.0053] & 14.2042 & $\underline{5.26 \times 10^{-5}}$ \\
 & ARIMA & -1.2679 [-1.2798, -1.2544] & 95.0155 & $\underline{< 10^{-100}}$ \\
 & Naive & -2.1110 [-2.1442, -2.0787] & 96.9455 & $\underline{< 10^{-100}}$ \\
\bottomrule
\end{tabular}
\end{table}


\begin{figure}[pos=!htbp,align=\centering]
    \centering
    \includegraphics[width=\linewidth]{figures/PIT_VAR_ES_calibration.png}
    \caption{Horizon-averaged calibration diagnostics for the probabilistic forecasting model. The PIT ECDF compares the empirical probability integral transform distribution to the uniform benchmark. VaR calibration reports empirical lower- and upper-tail exceedance rates against nominal tail levels. ES severity compares predicted expected shortfall with realized conditional tail means. Shaded regions denote pointwise 95\% moving-block bootstrap intervals over forecast origins.\href{https://github.com/QuantLet/onchain-insights/tree/main/12.\%20Evaluation\%20of\%20Probabilistic\%20Forecasts}{\includegraphics[height=0.6\baselineskip]{figures/Quantlet LOGO.png}}}
    \label{fig:calibration_diagnostics}
\end{figure}

\subsubsection{Backtesting and calibration diagnostics}

\begin{table}[pos=!htbp,align=\centering]
\centering
\small
\setlength{\tabcolsep}{6pt}
\begin{tabular}{l c c c c c c}
\toprule
Model & Horizon & $\alpha$ & VaR lower & VaR upper & ES lower & ES upper \\
\midrule
\multirow{4}{*}{Our Model} & \multirow{2}{*}{12h} & 0.5\% & \cellcolor{green!35}0.581 & \cellcolor{green!35}0.152 & \cellcolor{green!35}0.440 & \cellcolor{green!35}0.918 \\
&  & 1.0\% & \cellcolor{red!35}0.0136 & \cellcolor{green!35}0.872 & \cellcolor{red!35}0.042 & \cellcolor{green!35}0.659 \\
 & \multirow{2}{*}{24h} & 0.5\% & \cellcolor{green!35}0.802 & \cellcolor{green!35}0.336 & \cellcolor{green!35}0.698 & \cellcolor{green!35}0.296 \\
 &  & 1.0\% & \cellcolor{green!35}0.722 & \cellcolor{green!35}0.081 & \cellcolor{green!35}0.713 & \cellcolor{green!35}0.100 \\
\midrule
\multirow{4}{*}{GARCH}  & \multirow{2}{*}{12h}  & 0.5\% & \cellcolor{red!35}5.75e-4 & \cellcolor{red!35}7.85e-7 & \cellcolor{red!35}5e-4 & \cellcolor{red!35}5e-4 \\
 &  & 1.0\% & \cellcolor{red!35}0.014 & \cellcolor{red!35}9.32e-6 & \cellcolor{red!35}5e-4 & \cellcolor{red!35}5e-4 \\
 & \multirow{2}{*}{24h}  & 0.5\% & \cellcolor{red!35}0.001 & \cellcolor{red!35}6.63e-10 & \cellcolor{red!35}5e-4 & \cellcolor{red!35}5e-4 \\
 &  & 1.0\% & \cellcolor{red!35}1.72e-3 & \cellcolor{red!35}1.19e-5 & \cellcolor{red!35}5e-4 & \cellcolor{red!35}5e-4 \\
\midrule
\bottomrule
\end{tabular}
\caption{VaR unconditional-coverage (Kupiec) and expected-shortfall (Acerbi--Sz\'ekely $Z_2$) backtest p-values for lower and upper tails at the selected tail levels. Green cells indicate p-values $\geq 0.05$; red cells indicate rejection at 5\%. ES p-values are one-sided Monte Carlo values.}
\label{tab:var_es_backtests_comparison}
\end{table}



In addition to the CRPS, threshold-weighted CRPS, and tail quantile losses reported above, we assess calibration using PIT, VaR exceedance, and expected-shortfall diagnostics. The PIT, VaR-rate, and ES-severity plots summarize results across the 24 forecast horizons. Because forecast origins and horizons are serially dependent, uncertainty for these plots is estimated using a moving-block bootstrap over forecast origins. Each replication resamples contiguous 48-hour blocks and retains the complete 24-horizon forecast vector for each selected origin. We report pointwise 95\% percentile intervals.

Figure~\ref{fig:calibration_diagnostics} shows three complementary descriptive checks. The PIT mean is 0.5022, close to the uniform target of 0.5, and its variance is 0.0878, slightly above the uniform target of 0.0833. The PIT ECDF is close to the uniform benchmark ($D=0.0288$); a 48-hour block-bootstrap test does not reject uniformity at the 5\% level ($p=0.07896$, $B=2{,}000$), though the result is marginal at the 10\% level.

The VaR panel compares empirical exceedance rates with their nominal tail levels. At the larger levels (2--5\%), exceedance rates are above nominal in both tails, indicating under-coverage. Estimates at the smallest levels are less precise, as reflected in their wider bootstrap bands. Coverage appears closest to nominal around 0.5--1\%, levels relevant to rare-event monitoring.

The ES panel compares forecast ES with the realized conditional tail mean among observations exceeding the corresponding forecast VaR. Forecast ES is more extreme than the realized tail mean across the plotted levels. Since the bootstrap intervals overlap, particularly at the smallest levels, this pattern is consistent with conservative tail-severity forecasts but does not establish a statistically significant difference.

In addition to these descriptive diagnostics, we conduct VaR unconditional-coverage and ES backtests on non-overlapping forecast-origin sequences at the 12- and 24-hour horizons. Table~\ref{tab:var_es_backtests_comparison} reports Kupiec p-values for VaR and one-sided Monte Carlo p-values for the Acerbi--Székely $Z_2$ ES test at the 0.5\% and 1\% tail levels. For our model, the tests fail to reject at both levels and in both tails at 24 hours. At 12 hours, they fail to reject except for the lower tail at the 1\% level, where both VaR coverage ($p=0.0136$) and ES ($p=0.042$) are rejected. GARCH is rejected in every reported comparison. These results suggest better backtest performance for our model than for GARCH at the reported horizons and levels; they do not establish calibration at other horizons or tail levels. P-values are unadjusted across the reported tests.

\subsubsection{Interpretation of variable-selection weights in the forecasting architecture}

\begin{figure}[pos=!htbp,align=\centering]
\centering
\begin{subfigure}[b]{\textwidth}
    \includegraphics[width=\linewidth]{figures/variate_support_weight.png}
    \caption{Across-origin support frequency (upper panel), mean sparsemax weight, and mean weight conditional on nonzero support (lower panel). Variates are ordered by support frequency.}
    \label{fig:variate_support_weight}
\end{subfigure}

\begin{subfigure}[b]{\textwidth}
    \includegraphics[width=\linewidth]{figures/selection_raster.png}
    \caption{Binary selection raster over test origins. Blue indicates a nonzero sparsemax weight; the white background indicates zero weight. The red line reports the number of active variates at each origin.}
    \label{fig:variate_selection_raster}
\end{subfigure}

\caption{Variable-selection diagnostics from the forecasting model over the unseen test period. The GRN produces variable-selection scores that are transformed by sparsemax into nonnegative weights $\alpha_{t,j}$, with exact zeros for excluded variates. The weights describe the model's normalized allocation across variates in the variable-selection network; they are not probabilities of relevance or feature attributions.\href{https://github.com/QuantLet/onchain-insights/tree/main/13.\%20Sparse\%20Variate\%20Selection}{\includegraphics[height=0.6\baselineskip]{figures/Quantlet LOGO.png}}}
\label{fig:variate_selection}
\end{figure}

Figure~\ref{fig:variate_selection} shows a concentrated but origin-dependent variable-selection pattern. The degree-3 Gegenbauer 24-hour moving average is selected at virtually every test origin, while the second- and third-most frequent features are selected at roughly 70\% and 40\% of origins, respectively. The mean-weight panel shows a similarly concentrated allocation: the first feature receives an average weight of approximately 0.52, and the first two together receive about 70\% of the average selection weight.

The selected feature set is broader over the full test period than these leading frequencies alone suggest. In total, 68 variates receive nonzero sparsemax weight at least once, while 37 are selected at 1\% or more of test origins. The raster shows how this set changes over time: a small core recurs frequently, while additional features enter at particular origins. The number of active variates ranges from 1 to 14 per origin, with a mean of 5.15 and a standard deviation of 1.77. Thus, the entropy-regularized variable-selection network produces a sparse representation at each origin while retaining flexibility to use a wider pool of features across the test period.

Among the frequently selected inputs, Gegenbauer scores and energy features are prominent. In particular, the higher odd-degree Gegenbauer terms encode directional asymmetry in the liquidity curve around the peg and are informative about tail imbalance under the basis construction used here. The liquidation-count feature is also selected regularly, although its average selection weight is lower than that of the leading curve-shape features. Together, these patterns indicate that the forecasting model combines a recurring set of liquidity-curve signals with additional origin-specific features.

This selection profile notably differs from that of the early-warning classifier, whose attributions emphasize broader stress and market-state indicators. The forecasting model places more concentrated weight on specific liquidity-curve features, allowing its probabilistic forecast to condition on the local liquidity configuration at each origin. The comparison is qualitative: SHAP attributions and normalized sparsemax weights measure different quantities and their magnitudes are not directly comparable. The ablation results in the next subsection assess how these selected inputs impact forecast performance.

\subsubsection{Ablation study}

\begin{table}
\caption{Mean per-seed percentage change in forecast losses relative to the full model. Brackets report unadjusted 95\% paired Student-$t$ confidence intervals estimated from six seeds. Positive values indicate larger losses and therefore worse performance.}
\label{tab:ablation}
\resizebox{\textwidth}{!}{  
\begin{tabular}{lllllll}
\toprule
 Ablation & $\Delta\operatorname{CRPS} \;(\%)$ & $\Delta\operatorname{twCRPS} \;(\%)$ & $\Delta\operatorname{twCRPS_{low}} \;(\%)$ & $\Delta\operatorname{twCRPS_{high}} \;(\%)$ & $\Delta\operatorname{QL \,1\%} \;(\%)$ & $\Delta\operatorname{QL \,99\%} \;(\%)$\\
\midrule
No GPD tails & +0.5 [-0.0, +0.9]  & -0.2 [-2.4, +1.4]  & -0.5 [-2.8, +0.7]  & -0.1 [-2.9, +1.8]  & +57.0 [+36.1, +77.2]  & +58.4 [+37.0, +77.4]  \\
Gaussian output & +0.5 [+0.1, +1.0]  & +0.4 [-1.1, +1.7]  & +0.1 [-3.7, +2.6]  & +0.4 [-1.3, +2.0]  & +40.0 [+24.5, +55.3]  & +30.0 [+13.4, +43.1]  \\
No sparsity & +2.8 [+1.8, +3.9]  & +1.7 [-0.3, +3.3]  & +0.1 [-2.6, +3.4]  & +2.1 [-0.5, +4.0]  & +17.5 [+7.7, +28.8]  & +5.5 [-2.1, +13.5]  \\
No liquidity shape & +2.6 [+1.6, +3.7]  & +1.2 [-0.7, +2.9]  & +0.6 [-2.0, +4.4]  & +1.4 [-1.2, +3.4]  & +11.2 [+4.6, +19.0]  & +10.0 [+2.0, +16.2]  \\
Price only & +2.9 [+1.9, +4.0]  & +1.6 [-0.5, +3.2]  & +0.2 [-2.3, +3.3]  & +1.9 [-0.6, +4.1]  & +18.0 [+8.0, +29.1]  & +6.2 [-0.1, +12.1]  \\
\bottomrule
\end{tabular}}
\end{table}


Table~\ref{tab:ablation} reports the results of our ablation study, which evaluates the contribution of the probabilistic head and the on-chain features to forecast performance. For each metric, we report the percentage change relative to the full model, averaged across six paired random seeds, with unadjusted 95\% paired Student-$t$ confidence intervals.

The distributional ablations confirm the importance of the tail-focused probabilistic specification. Removing the GPD tails or using a Gaussian output head significantly increases both extreme quantile losses. The changes in aggregate scores are smaller: twCRPS intervals include zero for both ablations, while the Gaussian head also produces a small but positive change in CRPS.

The feature ablations produce moderate increases in aggregate losses, with CRPS increasing across all three specifications. Their mean twCRPS changes are also positive, although the confidence intervals include zero. Tail quantile losses generally increase as well, with the clearest effects at the 1\% quantile.

Threshold-weighted CRPS emphasizes the depeg-relevant region, but does not by itself show how model performance changes during depegs. We therefore compare ablation effects between calm and stressed regimes. A forecast horizon is classified as stressed if it crosses the 10~bps deviation threshold at least once. Figure~\ref{fig:stress_calm_ablation} shows that every ablation has larger losses than the full model in stressed periods. For the distributional ablations, their relatively favorable twCRPS changes in calm periods reverse to unfavorable changes under stress.

Overall, the results show that performance benefits from both the probabilistic specification and the sparse use of additional features. The model without the on-chain features remains a strong baseline, while the full model performs better in depeg-relevant periods and in the tails.

\begin{figure}[pos=!htbp,align=\centering]
    \centering
    \includegraphics[width=0.94\linewidth]{figures/stress_vs_calm.png}
    \caption{Mean percentage change in CRPS and twCRPS relative to the full model for calm and stressed forecast horizons. A horizon is classified as stressed if it crosses the 10~bps twCRPS threshold at least once.}
    \label{fig:stress_calm_ablation}
\end{figure}


\subsection{Economic evaluation}\label{sec:lp_results}

Table~\ref{tab:lp_main} reports the backtest over the holdout, and Figure~\ref{fig:lp_drawdowns} shows the
drawdowns over time. We discuss the benchmarks first, then the contribution of each model in turn.

\input{lp/tab_lp_main}

\begin{figure}[pos=!htbp,align=\centering]
    \centering
    \includegraphics[width=\linewidth]{figures/lp_drawdowns.png}
    \caption{Top: pool depeg, classifier alert episodes and 15\,bp events over the holdout. Bottom: drawdown
    of the net profit of the Gamma-style vault, the vault with forecaster ranges, and the two-model vault,
    which also pauses during the classifier's alerts (daily minimum).
    \href{https://github.com/QuantLet/onchain-insights/tree/main/14.\%20LP\%20vault\%20backtest}{\includegraphics[height=0.6\baselineskip]{figures/Quantlet LOGO.png}}}
    \label{fig:lp_drawdowns}
\end{figure}

\paragraph{Benchmarks.} The static range is almost riskless but earns only \lpStaticNet\,bp per year, because
the price spends much of its time away from parity. The Gamma-style vault keeps its capital in the active
range and earns \lpGammaNet\,bp per year, but at a cost: it repositions \lpGammaRepo{} times a month, its
LVR and costs absorb half of its fee income, and its risk is concentrated in the depeg events. Its maximum
drawdown (\lpGammaMaxDD\,bp) and worst day (\lpGammaWorstDay\,bp) both occur in the event of
4 February 2026, when the pool price reached \lpMaxDepeg\,bp.

\paragraph{The forecaster sizes the range.} Replacing the fixed $\pm2$-tick range with the forecaster's
interval raises net profit by \lpFcNetGain\,bp per year, to \lpFcNet\,bp, and the Sharpe ratio from
\lpGammaSharpe{} to \lpFcSharpe. The interval widens when the forecast uncertainty rises, which lowers the
vault's repositioning costs from \lpGammaCosts{} to \lpFcCosts\,bp per year; its worst day improves from \lpGammaWorstDay{} to \lpFcWorstDay\,bp and its daily CVaR from \lpGammaCVaR{}
to \lpFcCVaR\,bp. Figure~\ref{fig:lp_fees_lvr} decomposes these gains: within the event windows, the
forecaster range reduces LVR by \lpFcEvLVRCut\% and costs from \lpGammaEvCosts{} to \lpFcEvCosts\,bp while
keeping most of the fee income. Sizing alone does not protect against a sustained depeg, however: the range
is still repositioned as the price keeps moving, and the vault's maximum drawdown (\lpFcMaxDD\,bp) again
occurs in the largest event. Asymmetric intervals outperform symmetric ranges of the same width centred on
the median forecast with every forecaster (Table~\ref{tab:lp_forecasters}).

\begin{figure}[pos=!htbp,align=\centering]
    \centering
    \includegraphics[width=\linewidth]{figures/lp_fees_lvr.png}
    \caption{Fee income, LVR, repositioning costs and net profit of the three vaults within the event windows
    (24 hours before to 48 hours after each 15\,bp event) and outside them, summed over the holdout (bp of the
    vault's capital).
    \href{https://github.com/QuantLet/onchain-insights/tree/main/14.\%20LP\%20vault\%20backtest}{\includegraphics[height=0.6\baselineskip]{figures/Quantlet LOGO.png}}}
    \label{fig:lp_fees_lvr}
\end{figure}

\paragraph{The classifier protects against depegs.} Pausing the forecaster vault during the classifier's
alerts cuts its LVR within the event windows by \lpFcKEvLVRCut\% (from \lpFcEvLVR{} to \lpFcKEvLVR\,bp),
its event drawdown by \lpFcKEventDDCut\% (from \lpFcEventDD{} to \lpFcKEventDD\,bp) and its maximum
drawdown by \lpFcKMaxDDCut\% (from \lpFcMaxDD{} to \lpFcKMaxDD\,bp). The Sharpe ratio rises from
\lpFcSharpe{} to \lpFcKSharpe. Relative to the Gamma-style vault, the two-model vault reduces the maximum
drawdown by \lpFcKvsGammaMaxDDCut\% and the event drawdown by \lpFcKvsGammaEventDDCut\%, and it has the
smallest maximum drawdown and the highest Sharpe ratio of all actively managed vaults that do not use
future information.
The pause has the same effect on the Gamma-style vault, whose event drawdown falls by \lpGKEventDDCut\%,
from \lpGammaEventDD{} to \lpGKEventDD\,bp.

These reductions come from the timing of the alerts, not from pausing as such. Randomly timed pauses with
the classifier's episodes rarely achieve comparable reductions: the placebo $p$-values are \lpPFcEventDD{} for the event
drawdown and \lpPFcMaxDD{} for the maximum drawdown of the forecaster vault, and \lpPGammaEventDD{} and
\lpPGammaMaxDD{} for the Gamma-style vault (Figure~\ref{fig:lp_placebo} in Appendix~\ref{sec:lp_appendix}).
Lead time distinguishes the classifier from the price trigger, which acts on the same events but only once
the threshold is breached. With the classifier, the event drawdown of the forecaster vault is
\lpFcKEventDD\,bp, against \lpFcPEventDD\,bp with the price trigger and \lpFcOEventDD\,bp with the 24-hour
oracle, and it is smaller than with the price trigger in \lpFcKBetterEvents{} of the \lpNEvents{} events
(Table~\ref{tab:lp_events}). The ordering holds with every forecaster (Table~\ref{tab:lp_forecasters}) and in
every cost, size and competition scenario (Table~\ref{tab:lp_sensitivity}).

\input{lp/tab_lp_events}

\paragraph{Why the pause lowers net profit.} The pause is insurance, and it has a premium: it lowers net
profit by \lpFcKCost\,bp per year for the forecaster vault and \lpGKCost\,bp for the Gamma-style vault,
against \lpFcPCost{} and \lpGPCost\,bp for the price trigger. Figure~\ref{fig:lp_fees_lvr} shows where the
premium comes from. In this pool, depeg periods bring very large trading volumes, so a vault that keeps
repositioning earns fees that exceed its LVR on average: within the event windows the forecaster vault earns
\lpFcEvFees\,bp of fees against \lpFcEvLVR\,bp of LVR. A paused vault drifts out of range and forgoes both;
its fee income in the event windows falls to \lpFcKEvFees\,bp. In addition, \lpFalseEpisodes{} of the
classifier's \lpEpisodes{} episodes are not followed by an event and pause the vault in calm periods.
Figure~\ref{fig:lp_event_zoom} illustrates both cases. On 3 February 2025 the pause stops the accumulation of
LVR but also of fees. On 4 February 2026 the alerts start the day before the breach, and the paused vault
avoids the LVR that the other vaults incur when the price spikes to \lpMaxDepeg\,bp. The holdout contains no
severe depeg, so over this sample the premium exceeds the losses it avoids, as for any insurance whose value
lies in the tail. When half of the fee income is captured by just-in-time liquidity, the paused vaults earn
more than those that keep repositioning (Table~\ref{tab:lp_sensitivity}).

\begin{figure}[pos=!htbp,align=\centering]
    \centering
    \includegraphics[width=\linewidth]{figures/lp_event_zoom.png}
    \caption{Two holdout events. Top: pool depeg and classifier alert episodes; the red line marks the first
    15\,bp breach. Middle and bottom: cumulative fee income and cumulative LVR from the start of the window.
    The paused vault gives up fees and avoids LVR.
    \href{https://github.com/QuantLet/onchain-insights/tree/main/14.\%20LP\%20vault\%20backtest}{\includegraphics[height=0.6\baselineskip]{figures/Quantlet LOGO.png}}}
    \label{fig:lp_event_zoom}
\end{figure}

\paragraph{Out-of-sample check.} In 2024 the classifier's out-of-fold alerts detected \lpYDetected{} of the
\lpYNEvents{} events in the sample, so the forecaster's protection matters more (Table~\ref{tab:lp_2024}).
The event drawdown is \lpYGammaEventDD\,bp for the Gamma-style vault, \lpYGKEventDD\,bp with the classifier's
pauses, \lpYFcEventDD\,bp with the forecaster range and \lpYFcKEventDD\,bp with both. In this sample the
15\,bp price trigger performs similarly to the classifier once ranges are sized by the forecaster
(\lpYFcPEventDD\,bp), while for the Gamma-style vault the classifier gives the smaller maximum drawdown
(\lpYGKMaxDD{} against \lpYGPMaxDD\,bp).

\paragraph{A severe depeg.} Table~\ref{tab:lp_stress} and Figure~\ref{fig:lp_stress} illustrate the tail
that the pause insures against, using the USDC depeg of March 2023, when the pool price fell by up to
\lpStressMinDepeg\,bp. Fees are not observed in this period and the classifier's out-of-fold alerts do not
cover it, so the scenario reports losses only and starts the pause at, or a fixed number of hours before,
the first 15\,bp breach. Vaults that keep repositioning lose 20--30\% of their capital: the Gamma-style vault
repositions \lpStressGammaRepo{} times and loses \lpStressGamma\,bp, and the forecaster vault loses
\lpStressFc\,bp. Pausing at the breach limits the loss to a few basis points (\lpStressGammaBreach\,bp for
the Gamma-style vault), and pausing \lpStressLead{} hours earlier reduces it further
(\lpStressGammaLead\,bp). In such a crash the decisive protection is to stop repositioning around the event;
lead time adds a few basis points. In the moderate events of the holdout, where much of the damage occurs
before the breach, lead time halves the losses.

\input{lp/tab_lp_stress}

\begin{figure}[pos=!htbp,align=\centering]
    \centering
    \includegraphics[width=\linewidth]{figures/lp_stress_march2023.png}
    \caption{USDC depeg of March 2023. Top: pool depeg; the red line marks the first 15\,bp breach. Bottom:
    cumulative losses from LVR and repositioning costs, on a logarithmic scale; fee income is not observed in
    this period.
    \href{https://github.com/QuantLet/onchain-insights/tree/main/14.\%20LP\%20vault\%20backtest}{\includegraphics[height=0.6\baselineskip]{figures/Quantlet LOGO.png}}}
    \label{fig:lp_stress}
\end{figure}

\section{Conclusion and Perspectives} \label{sec:conclusion}

This paper develops a two-model framework for monitoring market-defined depeg risk in the USDC--USDT Uniswap V3 pool. The framework combines an early-warning classifier, which identifies the approach of a material price deviation, with a probabilistic forecaster, which characterizes the distribution and potential severity of future deviations. Both models draw on a broad set of on-chain and market information. A central contribution is the functional representation of the Uniswap V3 liquidity curve: liquidity depth, asymmetry, and local shape are encoded through economically interpretable features, including scores from a Gegenbauer polynomial basis. This representation captures information about executable liquidity that is not contained in the spot price or aggregate depth alone.

The out-of-sample results support the complementary roles of the two models. The selected early-warning policy detects 11 of 18 eligible events in development cross-validation and all five eligible events in the held-out period, where it generates approximately 2.5 false-alert episodes per month. Although the small number of held-out events limits inference about future detection rates, these results support the framework's potential for operational monitoring. The probabilistic model improves CRPS by 6.6\% and threshold-weighted CRPS by 14.2\% relative to the strongest benchmark for each metric, with statistically significant improvements against all five benchmarks. Reductions in extreme-quantile losses further indicate gains in forecasting the tails of the peg-deviation distribution. Together, the models provide complementary information about the occurrence, direction, and potential severity of prospective pool-price deviations.

An economic evaluation supports these complementary roles. In a backtest of liquidity provision to the
monitored pool, the forecaster sizes the vault's range: its predictive intervals raise net profit and the
Sharpe ratio relative to a Gamma-style vault, and reduce repositioning costs and daily tail losses. The
classifier protects against depegs: pausing repositioning during its alerts cuts the losses to arbitrageurs
around the held-out events by \lpFcKEvLVRCut\% and the drawdowns within them by \lpFcKEventDDCut\%, more
than a price trigger at the same threshold, which shows that its lead time has economic value. Together,
the two models reduce the vault's maximum drawdown by \lpFcKvsGammaMaxDDCut\%, at the cost of part of the
high fee income of depeg periods. In the March 2023 USDC depeg, vaults that keep repositioning lose 20--30\%
of their capital, the tail risk against which such protection is designed.

The feature-block ablations support the predictive relevance of liquidity conditions alongside price history and broader market information. For the early-warning classifier, removing the Uniswap V3 liquidity-curve feature block reduces detection from 11 to 4 of the 18 eligible cross-validation events, while increasing the false-alert burden. Removing Curve 3pool or Aave lending features also reduces event detection, suggesting complementary information across linked protocols. The probabilistic model repeatedly selects liquidity-curve scores, including odd-degree components that describe directional asymmetry around parity. Its ablations provide further evidence of their contribution: removing liquidity-shape features increases CRPS by 2.6\% and the 99th-quantile loss by 10.0\%. These findings support the predictive value of the liquidity representation, while leaving the economic mechanisms underlying these associations open to further investigation.

The empirical analysis is deliberately limited to venue-level deviations in a single pool. Its transferability to other markets therefore remains an empirical question. Applying the framework to additional Uniswap V3 pools and incorporating cross-chain activity, particularly on Ethereum Layer 2 venues, would broaden the assessment of liquidity conditions and test its robustness across trading environments. Higher-frequency observations could capture changes in pool state and order flow that hourly sampling may obscure. A further extension is to forecast prices and liquidity curves jointly. The compact Gegenbauer representation could support such forecasts and translate projected liquidity conditions into estimates of prospective trade slippage. The economic evaluation suggests further extensions: classifiers trained at the thresholds at which a given position is managed, and decisions that are costly to reverse within the hour, such as unwinding large or leveraged positions, where the lead time of the early-warning policy should be most valuable. These extensions would build on the study's central finding that the shape and distribution of executable liquidity provide useful information for monitoring and forecasting stablecoin price dislocations at the pool level.

\section{Acknowledgments}

This study has been conducted as part of the MSCA DIGITAL project. This project has received funding from the Horizon Europe research and innovation programme under the Marie Skłodowska-Curie Grant Agreement No. 101119635. The authors acknowledge the support of the Institute for Digital Assets (IDA), Bucharest University of Economic Studies, Romania.

\begin{figure}[pos=!ht, align=\centering]
    \includegraphics[width=0.21\linewidth]{figures/DIGITAL_logo.png}
    \includegraphics[width=0.2\linewidth]{figures/EU-logo.png}
\end{figure}

\section{Note on reproducibility}

To ensure transparency and reproducibility, all codes, trained models, and data preprocessing scripts used in this study are made publicly available as part of the Q2 (Quantinar-Quantlet) ecosystem.

Minimal working examples are provided in the Quantlet format, facilitating straightforward replication and extension of our results. The complete repository can be found at: \url{https://github.com/QuantLet/onchain-insights/tree/main}. The repository also contains a minimal working example of the pytorch repo used for model training. Finally, the forecast artifacts used to produce this paper's tables are available as a github release in this repo \url{https://github.com/QuantLet/onchain-insights/releases/tag/section12-forecasts-v1}. 

This commitment to open science is in alignment with the terms of the Marie Skłodowska-Curie Grant Agreement No. 101119635 and supports both the verification of our findings and the advancement of further research.


\printcredits
\bibliographystyle{cas-model2-names}

\bibliography{updated}

\newpage
\appendix

\section{Low rank functional analysis of the liquidity curve} \label{sec:liqcurve}

\subsection{Validity of low rank representation}

Our low-rank functional analysis of stablecoin liquidity provision goes as follows. We represent each observed liquidity curve as a functional object defined over a normalized tick grid. Specifically, for each observation time $t$, we consider the log-liquidity profile
\[
    \ell_t(x) = \log (1+L_t(x)),
\]
where $x$ denotes the relative tick index on a fixed window of $\pm 50$ ticks, and $L_t$ the active liquidity. Taking logarithms reduces the influence of extreme liquidity concentrations and allows the main cross-sectional structure of the curve to be captured more stably.

For stablecoin pools, two natural normalizations of the tick domain are possible. The first is a \emph{price-centered} representation (used in \cite{uniswapLPdynamics}), in which the curve is expressed relative to the current active tick. This representation captures the local liquidity environment faced by trades at the prevailing pool price. The second is a \emph{peg-centered} representation, in which the curve is expressed relative to the parity price of the stablecoin pair. This provides a fixed economic frame of reference and is particularly relevant for stablecoin pools, where liquidity providers may concentrate liquidity around the expected peg rather than around the transient market price. The temporal evolution of both representations is shown in Figure~\ref{fig:heatmaps}.

\begin{figure}[pos=!htbp,align=\centering]
    \centering
    \includegraphics[width=\textwidth]{figures/log_liquidity_heatmaps.png}
   \caption{Heatmap of Log-liquidity curves of the USDC/USDT Uniswap pool, centered around both the peg tick (parity) and the current active tick. The curves span a normalized range of $\pm 50$ ticks around the tick of interest. \href{https://github.com/QuantLet/onchain-insights/tree/main/5.\%20Functional\%20PCA\%20analysis\%20of\%20the\%20liquidity\%20curve}{\includegraphics[height=0.6\baselineskip]{figures/Quantlet LOGO.png}}
   }
    \label{fig:heatmaps}
\end{figure}

To assess whether these liquidity curves admit a parsimonious low-rank representation, we apply rolling functional principal component analysis (FPCA). For each rolling window, we compute the principal components of the log-liquidity curves and record the cumulative proportional variance explained (CPVE) by the first $K$ components, for $K=3,\ldots,7$. The resulting CPVE series are plotted in Figure~\ref{fig:pca_rolling}.


\begin{figure}[pos=!htbp,align=\centering]
    \centering
    \includegraphics[width=\linewidth]{figures/uniswap_liquidity_pca_analysis.png}
    \caption{Rolling FPCA of log-liquidity curves, both centered around the peg and around the current price. Principal components are estimated on a 1-month rolling window with a 24-hour stride. The cumulative proportion of variance explained (PVE) by the first 3 to 7 components is shown over time.
    \href{https://github.com/QuantLet/onchain-insights/tree/main/5.\%20Functional\%20PCA\%20analysis\%20of\%20the\%20liquidity\%20curve}{\includegraphics[height=0.6\baselineskip]{figures/Quantlet LOGO.png}}
    }    \label{fig:pca_rolling}
\end{figure}

The results indicate that both representations are well described by a relatively small number of functional components. In both the price-centered and peg-centered cases, the first seven components consistently explain more than 80\% of the variation in the liquidity curves. This supports the use of a low-rank approximation for the stablecoin liquidity surface and suggests that much of the observed variation in liquidity provision can be summarized by a small set of dominant deformation modes. The peg-centered representation exhibits a higher average CPVE, indicating that anchoring the curves at parity captures a particularly stable underlying structure. However, it also displays several sharper transient declines in explanatory power. By contrast, the price-centered representation appears somewhat more temporally stable in terms of CPVE, with fewer abrupt drops.

\subsection{Peg vs price-centered reference frame}

While CPVE establishes that a low-rank approximation is statistically adequate, interpretability is equally important for our study of depeg risk. In this respect, the peg-centered representation is more natural for stablecoin pools. Since parity is the fundamental reference price of the two assets, liquidity provision is expected to concentrate around the peg. Short-run deviations of the active pool price from parity are generally driven by trading pressure and pool microstructure, rather than by persistent changes in fundamental value. A price-centered representation therefore mechanically shifts the liquidity curve in response to these high-frequency price movements, potentially obscuring economically meaningful variation in the shape of the liquidity profile.

This effect is apparent in the temporal behaviour of low-order basis coefficients. Figure~\ref{fig:legendre_roughness} reports the first-degree Legendre score for both representations, together with a roughness measure based on normalized first differences. The price-centered scores fluctuate more rapidly, indicating that even coarse, low-frequency features of the curve vary noisily over time under the moving price frame. This temporal roughness hinders the interpretation of the corresponding low-rank factors, since changes in the scores may reflect re-centering around the active tick rather than genuine changes in liquidity provision. By contrast, the peg-centered scores evolve more smoothly and remain tied to a fixed economic reference point.

The same distinction is confirmed by the spatiotemporal spectral analysis in Figure~\ref{fig:spectrum}. In the peg-centered view, short-run price fluctuations appear mainly as localized, higher-spatial-frequency variation around parity. In the price-centered view, however, re-indexing the entire liquidity profile around the moving active tick introduces pronounced spectral intensity at low spatial frequencies and high temporal frequencies. This pattern is consistent with rapid variation in broad features of the curve, which is largely an artefact of the moving coordinate system rather than a stable liquidity-provision signal.

Consequently, although both representations admit a low-rank approximation, the peg-centered representation provides a more interpretable basis for studying depeg risk. Its fixed reference frame allows the leading functional components to be interpreted as economically meaningful deformations of liquidity around parity, such as changes in concentration at the peg, asymmetries between the two sides of the pool, or depletion of liquidity in the neighborhood of the peg. These modes are directly relevant to the resilience of the pool under depegging pressure and therefore form the primary representation used in the subsequent analysis.

\begin{figure}[pos=!htbp,align=\centering]
    \centering
    \includegraphics[width=\textwidth]{figures/spatiotemporal_spectrum_comparison.png}
   \caption{Spatio-temporal spectra of the peg- and price-centered log liquidity curves, as well as their log-power difference. The high temporal / low spatial frequency band is highlighted, showing much higher relative power in the price-centered representation. \href{https://github.com/QuantLet/onchain-insights/tree/main/5.\%20Functional\%20PCA\%20analysis\%20of\%20the\%20liquidity\%20curve}{\includegraphics[height=0.6\baselineskip]{figures/Quantlet LOGO.png}}
   }
    \label{fig:spectrum}
\end{figure}

\begin{figure}[pos=!htbp,align=\centering]
    \centering
    \includegraphics[width=\textwidth]{figures/coeff_roughness.png}
   \caption{Left: Temporal representation of the 1st degree score of the Legendre basis decomposition for both peg- and price-centered curves. Right: Temporal roughness of the Legendre scores for both curves, roughness is measured as the RMS first-difference of each score time series normalized by its standard deviation, so larger values indicate faster, noisier coefficient variation over time.\href{https://github.com/QuantLet/onchain-insights/tree/main/5.\%20Functional\%20PCA\%20analysis\%20of\%20the\%20liquidity\%20curve}{\includegraphics[height=0.6\baselineskip]{figures/Quantlet LOGO.png}}
   }
    \label{fig:legendre_roughness}
\end{figure}

\subsection{Gegenbauer bases and interpretable scores}

In order to construct an interpretable and temporally stable low-rank representation of the peg-centered log-liquidity curve, we use a fixed orthogonal polynomial basis. In contrast to FPCA eigenfunctions, which are estimated from the data and may change across rolling windows, a fixed basis yields scores that are directly comparable over time. This is particularly important in our depeg-risk setting, where the objective is not only to compress the liquidity curve, but also to track economically meaningful changes in its shape.

We decompose the peg-centered log-liquidity curve using a Gegenbauer basis. Gegenbauer polynomials generalize several standard orthogonal polynomial families, including Legendre and Chebyshev polynomials \citep{kde_cl}, through a parameter $\alpha$ that determines the weight function under which orthogonality is defined. For $\alpha>0$, the Gegenbauer polynomials $\{C_k^{(\alpha)}\}_{k\geq 0}$ are defined by the generating function
\begin{equation}
    \frac{1}{(1 - 2xz + z^2)^{\alpha}}
    =
    \sum_{k=0}^{\infty} C_k^{(\alpha)}(x) z^k,
    \qquad |z|<1,\quad x\in[-1,1].
\end{equation}
They are orthogonal on $[-1,1]$ with respect to the weight
\begin{equation}
    w_{\alpha}(x) = (1-x^2)^{\alpha-\frac{1}{2}},
\end{equation}
so that
\begin{equation}
    \int_{-1}^{1}
    C_j^{(\alpha)}(x) C_k^{(\alpha)}(x)
    w_{\alpha}(x)\,dx
    =
    h_k^{(\alpha)} \mathbf{1}_{\{j=k\}},
\end{equation}
where $h_k^{(\alpha)}$ denotes the squared norm of the $k$-th polynomial. The parameter $\alpha$ therefore controls the spatial emphasis of the approximation. Specifically in our use case, this parameter tunes the fit to the tail/center of the liquidity distribution. When $\alpha=1/2$, the weight is uniform and the basis reduces to the Legendre basis. Values $\alpha>1/2$ place more weight near the centre of the interval, while smaller values place relatively more emphasis on the tails.

Let $x\in[-1,1]$ denote the normalized peg-centered tick coordinate, with $x=0$ corresponding to parity. For each time $t$, we approximate the peg-centered log-liquidity curve $\ell_t(x)$ by the truncated expansion
\begin{equation}
    \tilde{\ell}_t(x)
    \approx
    \sum_{k=0}^{K} \beta_{t,k}\,\phi_k^{(\alpha)}(x),
\end{equation}
where $\phi_k^{(\alpha)}$ is the normalized Gegenbauer basis function of degree $k$, and $\beta_{t,k}$ is the corresponding time-varying score. In continuous form, these scores are obtained by projection,
\begin{equation}
    \beta_{t,k}
    =
    \int_{-1}^{1}
    \tilde{\ell}_t(x)\phi_k^{(\alpha)}(x) w_{\alpha}(x)\,dx,
\end{equation}
with the empirical implementation using the corresponding discrete weighted projection over the normalized tick grid.

\subsection{Selecting the Gegenbauer parameter $\alpha$}

The choice of $\alpha$ is not merely numerical; it affects both reconstruction accuracy and score interpretability. Figure~\ref{fig:alpha_impact} illustrates this trade-off. The left panel shows how increasing $\alpha$ concentrates the Gegenbauer weight around the peg, whereas lower values of $\alpha$ assign relatively more weight to the outer part of the tick range. This distinction matters for depeg risk. Under normal conditions, most liquidity is concentrated close to parity, but variation in this central region is not necessarily the most informative for depeg events. Instead, the relevant question is how much liquidity remains available as the price moves away from the peg. Tail regions of the liquidity curve are therefore economically important, because they describe the pool's ability to absorb larger dislocations.

The middle panel of Figure~\ref{fig:alpha_impact} shows the corresponding reconstruction trade-off for a fixed truncation level $K=7$. Larger values of $\alpha$ improve the fit near the centre of the interval, but this comes at the cost of a larger reconstruction error in the tails. Conversely, lower or moderate values of $\alpha$ preserve more information away from the peg. Since depeg episodes are precisely characterized by movement into these outer regions, a basis that over-prioritizes the centre may fit the usual near-peg liquidity profile well while discarding information that is most relevant under stress.

The right panel further shows that large values of $\alpha$ concentrate a substantial fraction of score variance in the degree-zero component. This weakens the interpretability of the decomposition, since the leading score then primarily captures shifts in the overall level of liquidity rather than meaningful changes in the shape of the curve. For depeg-risk analysis, however, higher-order shape features are essential: the first-degree score captures asymmetry across the peg, the second-degree score captures concentration or depletion around parity relative to the tails, and higher degrees capture more localized liquidity bands or gaps. Excessive concentration of variance in the degree-zero score therefore makes the representation less informative about the structural features of liquidity provision that determine resilience to price dislocations.

For this reason, we do not choose $\alpha$ solely to minimize central reconstruction error. Instead, we treat it as a hyperparameter governing the trade-off between smooth reconstruction, tail sensitivity, and score interpretability. Lower values of $\alpha$ retain more information in the tails, where liquidity conditions are most relevant for depeg risk, while avoiding excessive concentration of variance in the degree-zero component. The Gegenbauer scores then provide a compact and interpretable summary of the peg-centered curve: degree zero captures the overall liquidity level, degree one captures directional imbalance around the peg, degree two captures curvature and concentration, and higher degrees capture more localized shape changes. This intuition is supported by the cross-validation results for the early-warning models reported in Table~\ref{tab:cv_model_alpha_results}, where lower values of $\alpha$ yield better predictive performance, suggesting that they carry more economically meaningful information for depeg detection.

\begin{figure}[pos=!htbp,align=\centering]
    \centering
    \includegraphics[width=\textwidth]{figures/alpha_impact_decomposition.png}
   \caption{Effect of the Gegenbauer parameter $\alpha$ on the liquidity-curve representation. Left: spatial weighting induced by $w_\alpha(x)$ over the normalized peg-centered tick range, with larger $\alpha$ placing more weight near the peg. Middle: reconstruction trade-off for a fixed truncation $K=7$, showing centre, tail, and overall RMSE relative to the Legendre case $\alpha=0.5$. Right: distribution of score variance across polynomial degrees, illustrating that larger $\alpha$ concentrates variance in the degree-zero component and reduces the interpretability of shape-related scores.\href{https://github.com/QuantLet/onchain-insights/tree/main/5.\%20Functional\%20PCA\%20analysis\%20of\%20the\%20liquidity\%20curve}{\includegraphics[height=0.6\baselineskip]{figures/Quantlet LOGO.png}}
   }
    \label{fig:alpha_impact}
\end{figure}

\section{Early Warning policy Sensitivity Analysis} 
\label{sec:policy_sensitivity}
We examine sensitivity to the nominal hourly FPR target
$\eta$ used in Neyman--Pearson (NP) threshold calibration.
For comparison, the empirical policy selects the cutoff
that maximizes validation event recall subject to an
observed budget of five false-alert episodes per month.
All comparisons use identical chronological folds and a
12-hour quiet-period alert reset. NP calibration uses
24-hour negative spacing and nominal violation tolerance
$\delta_{\mathrm{NP}}=0.05$.

Table~\ref{tab:policy_sensitivity} reports the highest-recall
configuration under each policy. Every policy identifies a
foundation model, although the configuration selected at
$\eta=0.10$ is operationally inadmissible because its
false-alert upper confidence endpoint exceeds five episodes
per month. Among the NP configurations, observed hourly FPR
remains below the nominal target except at $\eta=0.10$.
These comparisons are development-CV results and do not
constitute independent validation of the selected policies.

The empirical policy and NP calibration at $\eta=0.08$
both detect 11/18 events. However, the empirical policy
produces a higher false-alert burden, at 3.53 rather than
3.11 episodes per month, and a higher observed hourly FPR.
We prioritize NP calibration because episode-based
constraints depend on the temporal clustering of threshold
exceedances: multiple false-positive hours can merge into
a single episode, allowing a low episode count despite
substantial hourly false-positive activity. NP calibration
provides an explicit hourly error target, although serial
dependence means that its formal iid guarantee is not
established here.

At $\eta=0.01$, three folds contain fewer than the 299
negative calibration observations required by the nominal
NP rule and therefore emit no alerts. The low recall at
this target consequently reflects calibration-data
limitations as well as the stringent FPR constraint.

\begin{table}[pos=!htbp,align=\centering]
    \centering
    \small
    \caption{Sensitivity to threshold-selection policy at the
    15-basis-point depeg threshold. Each row reports the
    highest-recall model--$\alpha$ configuration under that
    policy. False-alert intervals are calendar-block
    bootstrap 95\% confidence intervals. Infinite cutoffs
    impose a no-alert rule when NP calibration has insufficient
    negative observations. The dagger marks a configuration
    whose false-alert upper confidence endpoint exceeds
    five episodes per month.}
    \label{tab:policy_sensitivity}
    \resizebox{\textwidth}{!}{
    \begin{tabular}{llrrrrr}
        \toprule
        Policy & Model & $\alpha$ & Event recall & Observed hourly FPR & False alerts/month [95\% CI] & Folds with $c=\infty$ \\
        \midrule
        Empirical & Kumo & 1.0 & 11/18 (61.1\%) & 9.06\% & 3.53 [2.39, 4.87] & 0/4 \\
        NP $\eta=0.01$ & Kumo & 0.5& 1/18 (5.6\%)& 0.05\%& 0.08 [0.00, 0.26]& 3/4 \\
        NP $\eta=0.02$& TabICL-v2 & 0.5& 5/18 (27.8\%)& 1.01\%& 1.68 [0.83, 2.67] & 0/4 \\
        NP $\eta=0.05$& TabICL-v2 & 0.5& 8/18 (44.4\%)& 3.68\%& 2.52 [1.48, 3.73]& 0/4 \\
        NP $\eta=0.08$& Kumo & 1.0& 11/18 (61.1\%)& 7.38\%& 3.11 [1.90, 4.46]& 0/4 \\
        NP $\eta=0.10^{\dagger}$& Kumo & 0.5& 11/18 (61.1\%)& 13.95\%& 4.62 [3.28, 6.11]& 0/4 \\
        \bottomrule
    \end{tabular}
    }
\end{table}





\section{Probabilistic Forecasting Architecture} 
\label{sec:architecture_details}

\subsection{Stablecoin price residual analysis}\label{sec:residual_analysis}

Given the on-chain microstructure of stablecoin price formation, it is unlikely that price innovations are adequately described by simple parametric distributions such as the Gaussian. The time-varying liquidity distribution absorbs order flow in a highly nonlinear manner, producing very low volatility near the peg while still allowing for large price movements during periods of stress. To assess this empirically, we fit ARIMA models and examine whether the residuals are consistent with Gaussian and Student's $t$ distributions. As shown in Figure~\ref{fig:arima_qq}, both specifications provide a poor fit, indicating that standard innovation distributions, including heavy-tailed alternatives, do not adequately capture the stablecoin price regime. We further analyze residuals from GARCH models and obtain a similar result (Figure~\ref{fig:garch_qq}): even after accounting for heteroskedasticity and volatility clustering, residuals remain poorly described by standard parametric distributions. In particular, residual diagnostics reveal pronounced asymmetry between the lower and upper tails, suggesting that stablecoin price dynamics require a more flexible distributional forecasting framework.


\subsection{Implementation details}
\label{sec:hyperparameters}

Let $B$ denote batch size, $L$ input-window length, $V$ the number of
exogenous variables, $H$ forecast horizon, and $D$ hidden dimension. The
input has shape $B\times L\times(V+1)$; the final channel is the target
series and the preceding $V$ channels are the exogenous variables. The
reported run uses $L=168$, $V=98$, and $H=24$ (Table~\ref{tab:angel_hyperparameters}).

\paragraph{Normalization and target encoder}
Reversible instance normalization (RevIN) is applied separately to the
target and each exogenous channel, using statistics computed over the input
window. In the reported configuration, these are the window mean and
standard deviation, with learned affine parameters. The target's
normalization statistics are reused to transform the predicted location and
scale back to the target scale.

The normalized target is first projected to $C$ channels by a $1\times1$
convolution, then processed by $N$ causal residual blocks. Block $i$ uses
two causal convolutions with kernel size $k$ and dilation $2^i$, each
followed by GELU and dropout; the block output is passed through a residual
connection and layer normalization. The receptive field at the final
temporal position is
\begin{equation}
  1 + 2(k-1)\sum_{i=0}^{N-1}2^i
  = 1 + 2(k-1)(2^N-1).
\end{equation}
For the settings in Table~\ref{tab:angel_hyperparameters}, this is 125
time steps. The encoder concatenates the final temporal feature with the
mean of the temporal features, then projects the concatenation to $H$
horizon tokens of dimension $D$.

\paragraph{Covariate selection and fusion}
Each normalized exogenous trajectory is layer-normalized over time and
linearly projected from length $L$ to embedding dimension $D$. A learned variable embedding is added to each resulting token. The mean of the target horizon tokens is used as context for a gated residual network (GRN), which produces
one score per covariate. The scores are divided by a learned positive temperature and passed through sparsemax activation:
\begin{equation}
  w_{bv}
  = \operatorname{sparsemax}_{v}\!\left(
      \frac{g_{bv}}{\tau}
    \right),\qquad
  \tau=\operatorname{clip}\!\left(\exp(\gamma),\,0.1,\,10\right),
\end{equation}
where $g_{bv}$ is the context-conditioned score and $\gamma$ is a learned
scalar. The selected token is $w_{bv}$ times its covariate token; variables
with $w_{bv}\leq 10^{-8}$ are masked from attention. 

Two cross-attention layers use horizon embeddings as queries and selected covariate tokens as keys and values, with the settings listed in Table~\ref{tab:angel_hyperparameters}. The output head maps each fused horizon token to the distributional parameters. Attention diagnostics have one horizon-by-covariate map per layer; the reported diagnostic averages these maps over the fusion layers.

\paragraph{I-spline body and GPD tails}

For each horizon, the output head supplies a location, a scale, $K$ raw
spline coefficients, and two raw tail-shape parameters. The spline
coefficients are constrained to be positive,
\begin{equation}
  a_k = \operatorname{softplus}(\tilde a_k)+\epsilon,
  \qquad
  Q_0(u)=\sum_{k=1}^{K}a_k I_k(u),
  \qquad
  q_0(u)=\frac{dQ_0(u)}{du}
        =\sum_{k=1}^{K}a_k M_k(u),
\end{equation}
where $I_k$ and $M_k$ are the degree-three I-spline and M-spline bases.
This ensures that the body quantile function is nondecreasing. The basis
uses the knot placement and exponent listed in
Table~\ref{tab:angel_hyperparameters}.

The body is calibrated on the discrete quantile grid. Let $w_j$ be its
trapezoidal integration weights and define
\begin{equation}
  Q_c(u_j)=Q_0(u_j)-Q_0(0.5),\qquad
  \bar Q_c=\sum_j w_j Q_c(u_j),\qquad
  v_c=\sum_j w_j Q_c(u_j)^2-\bar Q_c^2.
\end{equation}
The calibrated body returned by the implementation is
\begin{equation}
  Q_B(u_j)=\ell_t+
        \frac{s_t}{\sqrt{v_c+\epsilon}}\,Q_c(u_j),
  \qquad
  q_B(u_j)=
        \frac{s_t}{\sqrt{v_c+\epsilon}}\,q_0(u_j),
\end{equation}
where $\ell_t$ and $s_t$ are the location and positive scale outputs after
the target's RevIN transformation. Thus $\ell_t$ anchors the body median, and $s_t$ sets its grid-based standard deviation up to numerical regularization.

The two GPD tails are joined to this body at $u_L$ and $u_U$. Their splice points and slope-matched scales are obtained by interpolation on the
quantile grid:
\begin{equation}
 z_L=Q_B(u_L),\quad z_U=Q_B(u_U),\qquad
 \beta_L=u_Lq_B(u_L),\quad
 \beta_U=(1-u_U)q_B(u_U).
\end{equation}
Each raw shape output $r_\xi$ is mapped to the configured range by
\begin{equation}
  \xi=\xi_{\min}+(\xi_{\max}-\xi_{\min})\operatorname{sigmoid}(r_\xi).
\end{equation}
For $u<u_L$ and $u>u_U$, the corresponding GPD quantile formulas in the
main text are used, with their continuous exponential limit when $\xi=0$.
The splice preserves the body quantile and its first derivative at each
join. The quantile grid is clamped away from 0 and 1 for numerical
stability, and its trapezoidal weights are used for both body calibration
and the quantile-based scores.

\paragraph{Training objective}
For each batch, the predictive loss is the mixture
\begin{equation}
  \mathcal{L}_{\mathrm{pred}}
  = \rho\,\mathcal{L}_{\mathrm{CRPS}}
    +(1-\rho)\,\mathcal{L}_{\mathrm{twCRPS}},
\end{equation}
using the quantile-grid approximations defined in the main text. The
selection penalty is the mean entropy of the sparsemax weights,
\begin{equation}
  \mathcal{H}
  =-\frac{1}{B}\sum_{b=1}^{B}\sum_{v=1}^{V}
       w_{bv}\log(w_{bv}+10^{-12}),
  \qquad
  \mathcal{L}_{\mathrm{total}}
  =\mathcal{L}_{\mathrm{pred}}+\lambda_0\mathcal{H}.
\end{equation}
The numerical settings for $\rho$, $\lambda_0$, the quantile grid, and
optimization are reported in Table~\ref{tab:angel_hyperparameters}.



\begin{table}[!htbp]
\centering
\small
\caption{Hyperparameters of the final distributional forecasting model. Settings are those of the run used for the paper results.}\label{tab:angel_hyperparameters}
\begin{threeparttable}
\begin{tabularx}{\linewidth}{@{}lYl@{}}
\toprule
\textbf{Component} & \textbf{Hyperparameter} & \textbf{Value} \\
\midrule
\multicolumn{3}{@{}l}{\textit{Forecasting task and input}} \\
& Input window & $168$ hours \\
& Forecast horizon & $24$ hours \\
& Historical exogenous features & $98$ \\
& Gegenbauer parameter, $\alpha$ & $0.1$ \\
\addlinespace
\multicolumn{3}{@{}l}{\textit{Distributional output and loss}} \\
& Quantile representation & I-spline, degree $3$ \\
& Knot placement / tail-density exponent & Power-tail / $3.0$ \\
& Number of basis functions & $5$ \\
& Tail model & generalized Pareto distribution (GPD) \\
& GPD splice quantiles $(u_{\mathrm{low}},u_{\mathrm{high}})$ & $(0.1,\,0.9)$ \\
& GPD shape range $[\xi_{\min},\xi_{\max}]$ & $[0,\,1]$ \\
& Quantile grid density & $256$ points \\
& Training loss & $\rho$CRPS + ($1-\rho$)twCRPS \\
& Loss mixture weight, $\rho$ & $0.7$ \\
& twCRPS thresholds / side / smoothing $h$ & $[-10,10]$ / two-sided / $1.0$ \\
\addlinespace
\multicolumn{3}{@{}l}{\textit{Network architecture}} \\
& Hidden dimension, $d_{\mathrm{model}}$ & $256$ \\
& Dropout & $0.0796$ \\
& RevIN & Enabled, affine \\
& Causal TCN channels / kernel size / levels & $128$ / $3$ / $5$ \\
& Covariate encoder layers / heads / feed-forward width & $2$ / $4$ / $256$ \\
& Variable-selector hidden width & $128$ \\
& Variable-selector mapping / initial temperature & Sparsemax / $2.0$ \\
& Cross-attention layers / heads / feed-forward width & $2$ / $2$ / $256$ \\
& Cross attention activation & Softmax \\
& Entropy sparsity penalty, $\lambda_{0}$ & $10^{-6}$ \\
\addlinespace
\multicolumn{3}{@{}l}{\textit{Optimization and training}} \\
& Training / test batch size & $20$ / $20$ \\
& Learning rate & $10^{-5}$ \\
& Maximum epochs / early-stopping patience & $100$ / $10$ \\
& Validation criterion & Loss \\
\addlinespace
\multicolumn{3}{@{}l}{\textit{Data split}} \\
& Validation split & $0.55$ \\
& Test split & $0.70$ \\
\bottomrule
\end{tabularx}
\end{threeparttable}
\end{table}

\section{Liquidity-provision backtest: additional results}\label{sec:lp_appendix}

Table~\ref{tab:lp_forecasters} applies the range rule to every forecaster, with asymmetric intervals (the main
specification) and with symmetric ranges of the same width centred on the median forecast.
Table~\ref{tab:lp_2024} repeats the comparison on 2024, Table~\ref{tab:lp_sensitivity} varies costs, vault
size, competition from just-in-time liquidity and fee reinvestment, and Figure~\ref{fig:lp_placebo} shows the
placebo distributions.

\input{lp/tab_lp_forecasters}

\input{lp/tab_lp_2024}

\input{lp/tab_lp_sensitivity}

\begin{figure}[pos=!htbp,align=\centering]
    \centering
    \includegraphics[width=\linewidth]{figures/lp_placebo.png}
    \caption{Event drawdown of the Gamma-style vault (left) and the forecaster vault (right) when the
    classifier's alert episodes are placed at random dates, compared with the vault without pauses, with the
    15\,bp price trigger, and with the classifier's alerts.
    \href{https://github.com/QuantLet/onchain-insights/tree/main/14.\%20LP\%20vault\%20backtest}{\includegraphics[height=0.6\baselineskip]{figures/Quantlet LOGO.png}}}
    \label{fig:lp_placebo}
\end{figure}

\section{Data Appendix}\label{sec:data-appendix}

\subsection{Data Sources}

\begin{table}[pos=!htbp,align=\centering]
\centering
\caption{The Graph subgraph identifiers}
\label{tab:subgraphs}
\resizebox{\textwidth}{!}{\begin{tabular}{lll}
\hline
\textbf{Protocol} & \textbf{Subgraph ID} & \textbf{Notes} \\
\hline
Uniswap~v3 (primary) & \texttt{5zvR82QoaXYFyDEKLZ9t6v9adgnptxYpKpSbxtgVENFV} & official Uniswap v3 \\
Uniswap~v3 (alt)     & \texttt{4cKy6QQMc5tpfdx8yxfYeb9TLZmgLQe44ddW1G7NwkA6} & fallback / cross-check \\
Curve (Messari)      & \texttt{3fy93eAT56UJsRCEht8iFhfi6wjHWXtZ9dnnbQmvFopF} & Messari standard schema \\
Chainlink            & \texttt{39QtNcs7YrvJUvh2sNVGFTKLeahMBpB3BKWypPrSNYLc} & USD index and FX oracles \\
Aave~v3 Ethereum     & \texttt{JCNWRypm7FYwV8fx5HhzZPSFaMxgkPuw4TnR3Gpi81zk} & \\
Aave~v2 Ethereum     & \texttt{C2zniPn45RnLDGzVeGZCx2Sw3GXrbc9gL4ZfL8B8Em2j} & Used until migration early 2023 \\
Ethereum Blocks           & \texttt{236pc6mnPH2EdGJxR5wunibyGsagq1JsSx5e2hx5tdoE} & Block timestamp indexing \\
\hline
\end{tabular}}
\end{table}

\begin{figure}[pos=!htbp,align=\centering]
    \centering
    \begin{subfigure}[b]{0.8\textwidth}
    \includegraphics[width=\linewidth]{figures/autoarima_qq_standardized.png}
    \caption{Standardized ARIMA residuals QQ-plots against Gaussian and Student's t distributions\href{https://github.com/QuantLet/onchain-insights/tree/main/9.\%20Parametric\%20quantile\%20function\%20characterisation}{\includegraphics[height=0.6\baselineskip]{figures/Quantlet LOGO.png}} \label{fig:arima_qq}}
    \end{subfigure}
    \begin{subfigure}[b]{0.8\textwidth}
    \includegraphics[width=\linewidth]{figures/garch_qq_standardized.png}
    \caption{Standardized GARCH(1,1) residuals QQ-plots against Gaussian and Student's t distributions \href{https://github.com/QuantLet/onchain-insights/tree/main/10.\%20Parametric\%20quantile\%20function\%20characterisation}{\includegraphics[height=0.6\baselineskip]{figures/Quantlet LOGO.png}}\label{fig:garch_qq}}
    \end{subfigure}
    \label{fig:residual_analysis}
\end{figure}
\begin{table}[pos=!htbp,align=\centering]
\centering
\caption{Main contract addresses (Ethereum Mainnet)}
\label{tab:pool_contracts}
\begin{tabular}{lll}
\hline
\textbf{Protocol} &\textbf{Contract} & \textbf{Address} \\
\hline
Curve & 3pool (AMM)  & \texttt{0xbebc44782c7db0a1a60cb6fe97d0b483032ff1c7} \\
Curve & 3CRV LP token & \texttt{0x6c3f90f043a72fa612cbac8115ee7e52bde6e490} \\
Curve & 3CRV Liquidity Gauge & \texttt{0xbFcF63294aD7105dEa65aA58F8AE5BE2D9d0952A} \\
\hline
Uniswap & USDC/USDT (0.01\%) & \texttt{0x3416cf6c708da44db2624d63ea0aaef7113527c6} \\
Uniswap & USDC/USDT (0.05\%) & \texttt{0x7858e59e0c01ea06df3af3d20ac7b0003275d4bf} \\
Uniswap & \texttt{NonFungiblePositionManager} & \texttt{0xc36442b4a4522e871399cd717abdd847ab11fe88} \\
\hline
AAVE v3 & USDT market & \texttt{0x23878914EFE38d27C4D67Ab83ed1b93A74D4086a} \\
AAVE v3 & USDC market & \texttt{0x98C23E9d8f34FEFb1B7BD6a91B7FF122F4e16F5c} \\
\hline
Chainlink & BTC/USD feed & \texttt{0xF4030086522a5bEEa4988F8cA5B36dbC97BeE88c} \\
Chainlink & ETH/USD feed & \texttt{0x5f4eC3Df9cbd43714FE2740f5E3616155c5b8419} \\
\hline
\end{tabular}
\end{table}

The data collection pipeline is fully reproducible. Collection scripts are publicly available in the project repository\footnote{\url{https://github.com/MSCA-DN-Digital-Finance/stablecoin-onchain-data}}. Data are collected nightly, and each release contains the most recent dataset in Parquet format.

The primary data source is The Graph. All subgraph queries are routed through The Graph's decentralised gateway:
\begin{center}
\url{https://gateway.thegraph.com/api/{API_KEY}/subgraphs/id/{SUBGRAPH_ID}}
\end{center}
Table~\ref{tab:subgraphs} provides a complete list of the subgraphs used.

To ensure temporal consistency across subgraphs, all observations are aligned to the Ethereum block closest to each hourly timestamp. Block timestamps and identifiers are obtained from the Ethereum Blocks subgraph. Historical position snapshots are queried at well-confirmed block heights, typically more than 128 blocks behind the chain head at the time of collection. Consequently, the probability that a chain reorganisation affects a stored observation is negligible.

The incremental updater records the last processed block number in a progress file. When resumed, it re-queries data beginning at that block. Any duplicate rows introduced by this overlap are removed through deduplication using the event primary key \texttt{(blockNumber, tx, logIndex, type, id)}.

Our secondary data source is Etherscan, which is accessed through its public V2 API (\texttt{https://api.etherscan.io/v2/api}, \texttt{chainid=1}). The platform is used to retrieve information that is unavailable from the selected subgraphs. In particular, event-level data on order flow for the 3pool LP token are collected from Etherscan. Incremental updates to these event-level data are likewise aligned by block number to the common hourly index.



\subsection{Full list of Dataset Features}

A complete list of features used in our dataset is provided in Table~\ref{tab:shap_feature_definitions}. Features are classified according to their economic category.
\clearpage
\begin{longtable}{p{0.24\textwidth} p{0.42\textwidth} p{0.24\textwidth}}
\caption{Feature definitions and economic group attribution.}
\label{tab:shap_feature_definitions} \\

\toprule
\textbf{Feature(s)} & \textbf{Short description} & \textbf{Economic group} \\
\midrule
\endhead

\bottomrule
\endfoot

\texttt{Gegenbauer\_*\_deg*},  
& Gegenbauer-basis coefficients describing the shape of the Uniswap V3 USDC/USDT liquidity curve. 
& Uniswap V3 liquidity-curve shape \\

\texttt{E\_*} 
& Energy measures derived from the Gegenbauer scores of the liquidity-curve. 
& Uniswap V3 liquidity-curve shape \\

\texttt{*\_ratio} 
& Ratios summarising relative energy based on frequency bands and parity. 
& Uniswap V3 liquidity-curve shape \\

\texttt{tangent\_up}, \texttt{tangent\_down} 
& Local slope of the swap-size price-impact curve for upward and downward trades. 
& Uniswap V3 liquidity-curve shape \\

\texttt{swap\_size\_imbalance} 
& Asymmetry between buy-side and sell-side swap-size capacity. 
& Uniswap V3 liquidity-curve shape \\

\texttt{tvlUSD\_100}, \texttt{tvlUSD\_500} 
& Total value locked in the Uniswap V3 USDC/USDT 1bp and 5bp fee-tier pools. 
& Uniswap V3 liquidity-curve shape \\

\midrule

\texttt{w\_USDC}, \texttt{w\_USDT} 
& USDC and USDT proportions in Curve 3pool. 
& Curve 3pool liquidity conditions \\

\texttt{curve\_entropy} 
& Entropy of Curve 3pool composition; lower values indicate greater imbalance. 
& Curve 3pool liquidity conditions \\

\texttt{gauge\_share\_3crv} 
& Share of 3pool LP tokens staked in the 3CRV gauge. 
& Curve 3pool liquidity conditions \\

\texttt{totalValueLockedUSD} 
& Total USD liquidity in Curve 3pool. 
& Curve 3pool liquidity conditions \\

\midrule

\texttt{eth\_*} 
& ETH market technical indicators and lagged market-state variables. 
& Broader market conditions \\

\texttt{btc\_*} 
& BTC market technical indicators and lagged market-state variables. 
& Broader market conditions \\

\texttt{usd\_index} 
& USD index, capturing broad dollar strength. 
& Broader market conditions \\

\texttt{fx\_volatility} 
& volatility of the USD index, capturing stability of the peg asset. 
& Broader market conditions \\

\texttt{fear\_greed\_index} 
& Crypto-market sentiment index. 
& Broader market conditions \\

\midrule

\texttt{depeg\_bps} 
& Current stablecoin peg deviation, measured in basis points. 
& Historical peg deviation \\

\texttt{depeg\_bps\_lag*} 
& Lagged peg deviations, capturing short-run persistence in depeg pressure. 
& Historical peg deviation \\

\midrule

\texttt{hhi\_24h\_rolling\_mean} 
& HHI of liquidity ownership concentration in NFPM position on Uniswap V3's USDC-USDT pool. 
& Liquidity ownership and position structure \\

\texttt{tick\_width\_24h\_rolling} 
& liquidity weighted average width of in-range Uniswap V3 liquidity positions over the previous 24 hours. 
& Liquidity ownership and position structure \\

\texttt{n\_in\_range\_log\_return} 
& Log return in the number of active in-range liquidity positions. 
& Liquidity ownership and position structure \\

\texttt{weighted\_mean\_age\_hours} 
& Liquidity weighted average age of in-range positions, in number of hours since mint.
& Liquidity ownership and position structure \\

\midrule

\texttt{swap\_count\_100}, \texttt{swap\_count\_500} 
& Number of swaps in the Uniswap V3 1bp and 5bp fee-tier pools. 
& Trading velocity and flows \\

\texttt{net\_amountUSD\_100}, \texttt{net\_amountUSD\_500} 
& Net USD-denominated swap flow in the Uniswap V3 1bp and 5bp pools. 
& Trading velocity and flows \\

\texttt{net\_amount\_token0}, \texttt{net\_amount\_token1} 
& Net swap flow into USDC (resp. USDT) in the Uniswap V3 1bp pool. 
& Trading velocity and flows \\

\texttt{hourlyVolumeUSD} 
& Hourly swap volume on the Curve 3pool. 
& Trading velocity and flows \\

\midrule

\texttt{supplied\_USD\_usdt}, \texttt{supplied\_USD\_usdc} 
& Total supplied USDT and USDC collateral on AAVE, measured in USD. 
& AAVE lending market conditions \\

\texttt{utilisation\_rate\_usdt}, \texttt{utilisation\_rate\_usdc} 
& Borrowing utilisation rates for USDT and USDC on AAVE. 
& AAVE lending market conditions \\

\texttt{liquidation\_USD} 
& Aggregate USD value of AAVE liquidations in previous hour. 
& AAVE lending market conditions \\

\end{longtable}



\end{document}